

\chapter{Anexos}


\section{Caracterización profesional}
\label{sec:CaractProf}

La caracterización profesional se aborda desde dos análisis complementarios, el primero de estos corresponde al análisis de las prácticas profesionales, y, el segundo, a la definición profesional. Para esto se considera cada análisis estructural de la actividad realizado con los docentes, según identidad matemática, y también una consulta con otros especialistas. 


\subsection{Análisis de las prácticas profesionales}


Este apartado se lleva a cabo por medio de categorías que se establecen para valorar las prácticas que, actualmente, trabajan los profesionales en Matemática, desde aquellas prácticas que han dejado de ser prioridades hasta las que son oportunidades para el matemático del futuro. Al respecto, se establecen las categorías de prácticas emergentes, dominantes y decadentes, con sus respectivas categorías y una forma de comprender el motivo por el cual han sido valoradas de esa forma.

\subsubsection{Prácticas emergentes}

\subsubsection*{Investigación fuera de la academia en disciplinas afines}


En el mundo, y, particularmente, en nuestro país, es notorio el crecimiento de la demanda de matemáticos en disciplinas tan variadas como las finanzas, la banca, la economía, la estadística. El trabajo interdisciplinario y las habilidades de liderazgo, para la conformación de equipos interdisciplinarios, se hacen cada vez más importantes e imprescindibles. 

El matemático de hoy en día debe estar en contacto permanente con los últimos acontecimientos científicos, no solo de la matemática propiamente, sino también de otras disciplinas afines. Debe saber entender y hacerse entender por otros profesionales, debe tener la empatía necesaria para interactuar, y dar respuestas satisfactorias a los requerimientos de otras disciplinas.

\subsubsection*{Análisis de datos}

La ciencia de los datos es aplicable, en términos generales, a diversas áreas (reconocimiento de imágenes, agrupamiento de datos, aprendizaje no supervisado, etc.). En resumen, la ciencia de
los datos está de moda. Las labores relacionadas con el manejo de datos a gran escala, ha hecho que el mercado, para los profesionales en áreas afines a la matemática, se expanda considerablemente. Esto, a su vez, ha provocado no solo una revalorización del matemático a nivel profesional, sino que, además, ha impactado la forma en que se hace mucha de la investigación académica, más de la mano del manejo e interpretación de datos. 

\subsubsection*{Experimentación científica cada vez más apoyada en la tecnología}

Los avances tecnológicos se han desencadenado de una manera impactante en las últimas décadas. Esto, sin duda, ha afectado de manera muy positiva el desarrollo de la investigación en matemática. Por citar ejemplos, la simulación de fenómenos naturales, que tiempo atrás se realizaba más de la mano de simulaciones mecánicas o mediante la observación directa, hoy en día es posible realizarla de manera más versátil mediante el uso de la tecnología. Al estudiar fenómenos financieros o económicos, o, en general, fenómenos de comportamiento social, las simulaciones estocásticas son una herramienta que, si bien han existido por décadas, cada vez están más a la mano y de manera más natural. El impacto de estos recursos para la investigación es innegable. 

\subsubsection{Prácticas dominantes}

\subsubsection*{Investigación académica}

La investigación académica, tanto teórica como aplicada, ha estado presente como una práctica profesional en matemática, posiblemente desde la antigüedad. Esta es una práctica que definitivamente no va a pasar de moda. El matemático es posiblemente uno de los científicos al que se le facilita más la investigación, sobre todo si esta es de carácter teórico, debido a que no requiere de realización de grandes experimentos en laboratorios sofisticados. Muchos grandes avances se pueden realizar solo con disciplina, lápiz y papel. 


\subsubsection*{Docencia Universitaria }


Es muy difícil desligar a un matemático de la Docencia Universitaria. Si bien, algunas prácticas emergentes han provocado que el sector productivo absorba ciertos profesionales del área de la matemática, una gran mayoría de matemáticos en el mundo mantienen su ligamen con la academia. Los grandes institutos de investigación en matemática alrededor del mundo que no pertenecen a las Universidades mantienen estrecha relación con estas.

\subsubsection*{Modelización de fenómenos de otras disciplinas}

Un área importante de la matemática, posiblemente la que más se reconoce desde fuera de la matemática, es la de modelización, o, en general, la matemática aplicada. Si bien hay varias prácticas emergentes relacionadas con esta área de la matemática, lo cierto es que, globalmente, representa una práctica que ha sido dominante por mucho tiempo. En tiempos recientes, esta práctica se ha venido desarrollando en diferentes áreas donde los matemáticos juegan un papel importante en el desarrollo de nuevas herramientas de modelización. 

\subsubsection*{Trabajo interdisciplinario}


Si bien algunas prácticas emergentes están relacionadas de cerca con el trabajo interdisciplinario, esta práctica ha estado presente por largo tiempo, por ejemplo, en aplicaciones en Ingeniería y Física, en donde la resolución numérica de ecuaciones diferenciales parciales ayuda a modelar fenómenos en diversas aplicaciones (electromagnetismo, fluidos, etc.). Otro ejemplo importante de trabajo interdisciplinario es el modelado de enfermedades infecciosas, con sus posibles impactos en política de salud pública. 

El objetivo de un trabajo interdisciplinario consiste en abordar de forma sistemática una pregunta de investigación, por medio de la articulación de diversas áreas de conocimiento, a fin de que sus actividades no se produzcan en forma aislada, dispersa o fraccionada. Este trabajo integrado potencializa el conocimiento teórico y metodológico necesario para abordar los complejos y diversos retos que enfrenta la ciencia del siglo XXI. Los matemáticos juegan un papel fundamental en el trabajo interdisciplinario, donde sirven como traductores de un problema de la vida real al lenguaje matemático. Esto permite utilizar las herramientas matemáticas disponibles. 

\subsubsection{Decadentes}

\subsubsection{Labores de índole actuarial}

Tal vez, esta sea ya una práctica que desapareció por completo de la matemática. Al menos en nuestro país, con una carrera de Ciencias Actuariales relativamente joven, las labores de tipo actuarial como los peritajes legales, asesorías en seguros o pensiones y labores por el estilo, son ahora asumidas principalmente por los actuarios, por lo que dejan de ser una opción para los matemáticos.


\subsubsection{Docencia en enseñanza secundaria}

Debido a la existencia, cada vez mayor, de profesionales especializados en el área de Educación Matemática, es notoria la disminución de la oferta laboral para los matemáticos de parte de este sector, al menos en nuestro país.  

De esta forma, se presenta un profesional en matemática que está en constante cambio. Hay necesidad de dar respuesta a las demandas actuales de una sociedad que requiere profesionales en matemática que no solo estén presentes en un entorno académico, sino también en los entornos productivos de la industria y el comercio. Se han ampliado los escenarios donde la matemática participa, y, por tanto, en los que el matemático puede ingresar para aprender, descubrir, investigar y, sobre todo, vivir con pasión.

\subsection{Definición profesional}

A partir de las identidades matemáticas designadas en el Análisis de la Disciplina, se estableció un objeto ligado a la identidad matemática que se caracterizó desde el punto de vista profesional de la siguiente forma:

%\rowcolor[HTML]{00C0F3} 
%\multicolumn{3}{|c|}{\cellcolor[HTML]{00C0F3}Reuniones de comité} \\ \hline
%\rowcolor[HTML]{8ED8F8}

\begin{center}
    \textbf{Objeto 1: Estructuras asociadas a clases de objetos y relaciones entre ellas.}
\end{center}

\begin{table}[H]
\begin{tabular}{|cl|}
\hline
\rowcolor[HTML]{00C0F3} 
\multicolumn{2}{|c|}{\cellcolor[HTML]{00C0F3}} \\
\rowcolor[HTML]{00C0F3} 
\multicolumn{2}{|c|}{\cellcolor[HTML]{00C0F3}} \\
\rowcolor[HTML]{00C0F3} 
\multicolumn{2}{|c|}{\multirow{-3}{*}{\cellcolor[HTML]{00C0F3}\textbf{\begin{tabular}[c]{@{}c@{}}Objetivo 1.1: Clasificar estructuras algebraicas de acuerdo con\\ los morfismos que las preservan para reconocerlas y relacionarlas.\\ .\end{tabular}}}} \\ \hline
\rowcolor[HTML]{8ED8F8} 
\multicolumn{1}{|c|}{\cellcolor[HTML]{8ED8F8}\textbf{\begin{tabular}[c]{@{}c@{}}.\\ Acciones\\ .\end{tabular}}} & \multicolumn{1}{c|}{\cellcolor[HTML]{8ED8F8}\textbf{\begin{tabular}[c]{@{}c@{}}.\\ Estrategias metodológicas \\ y tecnológicas\\ .\end{tabular}}} \\ \hline
\multicolumn{1}{|l|}{\begin{tabular}[c]{@{}l@{}}.\\ Comprender estructuras algebraicas básicas y \\ concretas (números reales y complejos, \\ polinomios y soluciones).\\ .\end{tabular}} & \begin{tabular}[c]{@{}l@{}}.\\ Exploración bibliográfica y web\\ .\end{tabular} \\ \hline
\multicolumn{1}{|l|}{\begin{tabular}[c]{@{}l@{}}.\\ Comprender las estructuras algebraicas (conjuntos\\ y relaciones, espacios vectoriales, grupos, anillos, \\ cuerpos, módulos,categorías...).\\ .\end{tabular}} & \begin{tabular}[c]{@{}l@{}}.\\ Resolución de ejercicios.\\ .\end{tabular} \\ \hline
\multicolumn{1}{|l|}{\begin{tabular}[c]{@{}l@{}}.\\ Estudiar interacciones entre estructuras de distinto tipo\\ .\end{tabular}} & \begin{tabular}[c]{@{}l@{}}.\\ Trabajo en equipo\\ .\end{tabular} \\ \hline
\end{tabular}
\end{table}

\begin{table}[H]
\begin{tabular}{|cl|}
\hline
\rowcolor[HTML]{00C0F3} 
\multicolumn{2}{|c|}{\cellcolor[HTML]{00C0F3}} \\
\rowcolor[HTML]{00C0F3} 
\multicolumn{2}{|c|}{\cellcolor[HTML]{00C0F3}} \\
\rowcolor[HTML]{00C0F3} 
\multicolumn{2}{|c|}{\multirow{-3}{*}{\cellcolor[HTML]{00C0F3}\textbf{\begin{tabular}[c]{@{}c@{}}Objetivo 1.2: Comprender la interacción del Álgebra con otras áreas del conocimiento\\ para utilizar métodos algebraicos en la resolución de problemas de otras áreas.\\ \end{tabular}}}} \\ \hline
\rowcolor[HTML]{8ED8F8} 
\multicolumn{1}{|c|}{\cellcolor[HTML]{8ED8F8}\textbf{\begin{tabular}[c]{@{}c@{}}\\ Acciones\\ \end{tabular}}} & \multicolumn{1}{c|}{\cellcolor[HTML]{8ED8F8}\textbf{\begin{tabular}[c]{@{}c@{}}\\ Estrategias metodológicas \\ y tecnológicas\\ \end{tabular}}} \\ \hline
\multicolumn{1}{|l|}{\begin{tabular}[c]{@{}l@{}}\\ Estudiar las áreas relacionadas (geometría algebraica, \\ teoría de números, topología, análisis funcional, \\ lógica).\\ \end{tabular}} & \begin{tabular}[c]{@{}l@{}} \\ Exploración bibliográfica.\\ \end{tabular} \\ \hline
\multicolumn{1}{|l|}{\begin{tabular}[c]{@{}l@{}}\\ Estudiar las aplicaciones de resultados algebraicos \\ en otras áreas.\\ \end{tabular}} & \begin{tabular}[c]{@{}l@{}}\\ Interacción con expertos de otras \\ áreas.\\ \end{tabular} \\ \hline
\multicolumn{1}{|l|}{\begin{tabular}[c]{@{}l@{}}\\ Estudiar las aplicaciones de otras áreas al álgebra\\ \end{tabular}} & \begin{tabular}[c]{@{}l@{}}\\ Desarrollo de proyectos y\\  presentaciones.\\ \end{tabular} \\ \hline
\rowcolor[HTML]{00C0F3} 
\multicolumn{2}{|c|}{\cellcolor[HTML]{00C0F3}} \\
\rowcolor[HTML]{00C0F3} 
\multicolumn{2}{|c|}{\cellcolor[HTML]{00C0F3}} \\
\rowcolor[HTML]{00C0F3} 
\multicolumn{2}{|c|}{\multirow{-3}{*}{\cellcolor[HTML]{00C0F3}\textbf{\begin{tabular}[c]{@{}c@{}}Objetivo 1.3: Contribuir a la creación de conocimiento desde el Álgebra \\ para proveer nuevos métodos algebraicos en la matemática o fuera de ella.\\ \end{tabular}}}} \\ \hline
\rowcolor[HTML]{8ED8F8} 
\multicolumn{1}{|c|}{\cellcolor[HTML]{8ED8F8}\textbf{\begin{tabular}[c]{@{}c@{}}\\ Acciones\\ \end{tabular}}} & \multicolumn{1}{c|}{\cellcolor[HTML]{8ED8F8}\textbf{\begin{tabular}[c]{@{}c@{}}\\ Estrategias metodológicas \\ y tecnológicas\\ \end{tabular}}} \\ \hline
\multicolumn{1}{|l|}{\begin{tabular}[c]{@{}l@{}}\\ Crear nuevas estructuras algebraicas y nuevas \\ técnicas, como respuesta a problemas del Álgebra\\  y de otras áreas.\\ \end{tabular}} & \begin{tabular}[c]{@{}l@{}}\\ Exploración bibliográfica\\ \\ \end{tabular} \\ \hline
\multicolumn{1}{|l|}{\begin{tabular}[c]{@{}l@{}} \\ Aplicar técnicas algebraicas a resolución de nuevos\\ problemas.\\ \end{tabular}} & \begin{tabular}[c]{@{}l@{}}\\ Interacción con expertos de otras\\  áreas.\\ \end{tabular} \\ \hline
\multicolumn{1}{|l|}{\begin{tabular}[c]{@{}l@{}} \\ Usar conocimiento existente en la resolución de \\ problemas  en Álgebra.\\  \end{tabular}} & \begin{tabular}[c]{@{}l@{}} \\ Seminario, charlas. \\  \end{tabular} \\ \hline
\rowcolor[HTML]{00C0F3} 
\multicolumn{2}{|c|}{\cellcolor[HTML]{00C0F3}} \\
\rowcolor[HTML]{00C0F3} 
\multicolumn{2}{|c|}{\cellcolor[HTML]{00C0F3}} \\
\rowcolor[HTML]{00C0F3}
\multicolumn{2}{|c|}{\multirow{-3}{*}{\cellcolor[HTML]{00C0F3}\textbf{\begin{tabular}[c]{@{}c@{}}Objetivo 1.4: Divulgar la producción de conocimientos, \\ para que pueda ser comprendida y usada\\ \end{tabular}}}} \\ \hline
\rowcolor[HTML]{8ED8F8} 
\multicolumn{1}{|c|}{\cellcolor[HTML]{8ED8F8}\textbf{\begin{tabular}[c]{@{}c@{}} \\ Acciones\\ \end{tabular}}} & \multicolumn{1}{c|}{\cellcolor[HTML]{8ED8F8}\textbf{\begin{tabular}[c]{@{}c@{}}\\ Estrategias metodológicas\\  y tecnológicas\\ \end{tabular}}} \\ \hline
\multicolumn{1}{|l|}{\begin{tabular}[c]{@{}l@{}}\\ Elaborar documentos formales\\ \end{tabular}} & \begin{tabular}[c]{@{}l@{}}\\ Uso de sofware especializado para \\ redactar artículos científicos.\\ \end{tabular} \\ \hline
\multicolumn{1}{|l|}{\begin{tabular}[c]{@{}l@{}}\\ Socializar los conocimientos mediante la \\ interacción con otros profesionales\\ \end{tabular}} & \begin{tabular}[c]{@{}l@{}}\\ Asistencia a conferencias, congresos, \\ seminarios, etc.\\ \end{tabular} \\ \hline
\multicolumn{1}{|l|}{\begin{tabular}[c]{@{}l@{}}\\ Comunicar los resultados mediante un lenguaje\\ interdisciplinario\\ \end{tabular}} & \begin{tabular}[c]{@{}l@{}}\\ Realizar informes y presentaciones.\\ \end{tabular} \\ \hline
\end{tabular}
\end{table}

\vspace{2mm}

\begin{center}
    \textbf{Objeto 2: Propiedades de las estructuras geométricas}
\end{center}

\begin{table}[H]
\begin{tabular}{|ll|}
\hline
\rowcolor[HTML]{00C0F3} 
\multicolumn{2}{|c|}{\cellcolor[HTML]{00C0F3}} \\
\rowcolor[HTML]{00C0F3} 
\multicolumn{2}{|c|}{\cellcolor[HTML]{00C0F3}} \\
\rowcolor[HTML]{00C0F3} 
\multicolumn{2}{|c|}{\multirow{-3}{*}{\cellcolor[HTML]{00C0F3}\textbf{\begin{tabular}[c]{@{}c@{}}Objetivo 2.1: Conocer los principales tipos de Geometría para visualizar\\ los espacios euclídeos, hiperbólicos, esféricos y proyectivos.\end{tabular}}}} \\ \hline
\rowcolor[HTML]{8ED8F8} 
\multicolumn{1}{|c|}{\cellcolor[HTML]{8ED8F8}\textbf{Acciones}} & \multicolumn{1}{c|}{\cellcolor[HTML]{8ED8F8}\textbf{\begin{tabular}[c]{@{}c@{}}Estrategias metodológicas \\ \\ y tecnológicas\end{tabular}}} \\ \hline
\multicolumn{1}{|l|}{\begin{tabular}[c]{@{}l@{}}Estudiar las bases de la geometría clásica\\ (euclidiana).\end{tabular}} &  \\ \hline
\multicolumn{1}{|l|}{\begin{tabular}[c]{@{}l@{}}Estudiar las bases de Álgebra y Análisis real\\ y complejo\end{tabular}} & Exploración bibliográfica y web \\ \hline
\multicolumn{1}{|l|}{Estudiar la geometría proyectiva real y compleja.} & Resolución de ejercicios \\ \hline
\multicolumn{1}{|l|}{Reconocer los diferentes tipos de geometría.} & Trabajo en equipo \\ \hline
\multicolumn{1}{|l|}{Enfocar una o algunas ramas de preferencia.} &  \\ \hline
\rowcolor[HTML]{00C0F3} 
\multicolumn{2}{|c|}{\cellcolor[HTML]{00C0F3}} \\
\rowcolor[HTML]{00C0F3} 
\multicolumn{2}{|c|}{\cellcolor[HTML]{00C0F3}} \\
\rowcolor[HTML]{00C0F3} 
\multicolumn{2}{|c|}{\multirow{-3}{*}{\cellcolor[HTML]{00C0F3}\textbf{\begin{tabular}[c]{@{}c@{}}Objetivo 2.2: Conocer las propiedades de curvas, superficies y variedades para visualizar \\ las nociones geométricas básicas en espacios de dos, tres y más dimensiones.\\ .\end{tabular}}}} \\ \hline
\rowcolor[HTML]{8ED8F8} 
\multicolumn{1}{|c|}{\cellcolor[HTML]{8ED8F8}\textbf{Acciones}} & \multicolumn{1}{c|}{\cellcolor[HTML]{8ED8F8}\textbf{\begin{tabular}[c]{@{}c@{}}Estrategias metodológicas \\ \\ y tecnológicas\end{tabular}}} \\ \hline
\multicolumn{1}{|l|}{Estudiar las propiedades conocidas} &  \\ \hline
\multicolumn{1}{|l|}{\begin{tabular}[c]{@{}l@{}}Estudiar objetos asociados a las estructuras \\ geométricas (haces, fibrados vectoriales, \\ curvaturas).\end{tabular}} & Resolución de ejercicios \\ \hline
\multicolumn{1}{|l|}{\begin{tabular}[c]{@{}l@{}}Indagar los métodos y técnicas de estudio para \\ resolver los problemas geométricos.\end{tabular}} & \begin{tabular}[c]{@{}l@{}}Interacción con expertos de otras\\ áreas\end{tabular} \\ \hline
\rowcolor[HTML]{00C0F3} 
\multicolumn{2}{|c|}{\cellcolor[HTML]{00C0F3}} \\
\rowcolor[HTML]{00C0F3} 
\multicolumn{2}{|c|}{\cellcolor[HTML]{00C0F3}} \\
\rowcolor[HTML]{00C0F3} 
\multicolumn{2}{|c|}{\multirow{-3}{*}{\cellcolor[HTML]{00C0F3}\textbf{\begin{tabular}[c]{@{}c@{}}Objetivo 2.3: Resolver problemas de las estructuras geométricas\\ \\ para producir nuevo conocimiento.\end{tabular}}}} \\ \hline
\rowcolor[HTML]{8ED8F8} 
\multicolumn{1}{|c|}{\cellcolor[HTML]{8ED8F8}\textbf{\begin{tabular}[c]{@{}c@{}}.\\ Acciones\\ .\end{tabular}}} & \multicolumn{1}{c|}{\cellcolor[HTML]{8ED8F8}\textbf{\begin{tabular}[c]{@{}c@{}}.\\ Estrategias metodológicas \\ y tecnológicas\\ .\end{tabular}}} \\ \hline
\multicolumn{1}{|l|}{Reconocer el problema por resolver} & Exploración bibliográfica \\ \hline
\multicolumn{1}{|l|}{\begin{tabular}[c]{@{}l@{}}Identificar aplicaciones a otros problemas en \\ Geometría, otras áreas de la Matemática o en \\ la vida real.\end{tabular}} & \begin{tabular}[c]{@{}l@{}}Interacción con expertos de otras\\ \\  áreas.\end{tabular} \\ \hline
\multicolumn{1}{|l|}{\begin{tabular}[c]{@{}l@{}}Entender los pasos sucesivos al problema que\\ se trata de resolver.\end{tabular}} & Seminario, charlas. \\ \hline
\multicolumn{1}{|l|}{\begin{tabular}[c]{@{}l@{}}Plantear las técnicas necesarias para resolver el\\  problema\end{tabular}} &  \\ \hline
\end{tabular}
\end{table}


\begin{table}[H]
\begin{tabular}{|ll|}
\hline
\rowcolor[HTML]{00C0F3} 
\multicolumn{2}{|c|}{\cellcolor[HTML]{00C0F3}} \\
\rowcolor[HTML]{00C0F3} 
\multicolumn{2}{|c|}{\cellcolor[HTML]{00C0F3}} \\
\rowcolor[HTML]{00C0F3} 
\multicolumn{2}{|c|}{\multirow{-3}{*}{\cellcolor[HTML]{00C0F3}\textbf{\begin{tabular}[c]{@{}c@{}}Objetivo 2.4: Investigar problemas en geometría para desarrollar\\ aplicaciones a problemas reales.\end{tabular}}}} \\ \hline
\rowcolor[HTML]{8ED8F8} 
\multicolumn{1}{|c|}{\cellcolor[HTML]{8ED8F8}\textbf{Acciones}} & \multicolumn{1}{c|}{\cellcolor[HTML]{8ED8F8}\textbf{\begin{tabular}[c]{@{}c@{}}Estrategias metodológicas\\ \\  y tecnológicas\end{tabular}}} \\ \hline
\multicolumn{1}{|l|}{\begin{tabular}[c]{@{}l@{}}Estudiar modelos geométricos de problemas \\ reales abiertos.\end{tabular}} &  \\ \hline
\multicolumn{1}{|l|}{\begin{tabular}[c]{@{}l@{}}Interactuar con profesionales de otras \\ disciplinas.\end{tabular}} & Visualización. \\ \hline
\multicolumn{1}{|l|}{Encriptar la información.} & Sofware especializado. \\ \hline
\multicolumn{1}{|l|}{Simplificar el problema para que sea aplicable.} & Interacción con otras disciplinas. \\ \hline
\multicolumn{1}{|l|}{\begin{tabular}[c]{@{}l@{}}Describir el problema mediante números que \\ puedan ser insertados en un sistema.\end{tabular}} &  \\ \hline
\rowcolor[HTML]{00C0F3} 
\multicolumn{2}{|c|}{\cellcolor[HTML]{00C0F3}} \\
\rowcolor[HTML]{00C0F3} 
\multicolumn{2}{|c|}{\cellcolor[HTML]{00C0F3}} \\
\rowcolor[HTML]{00C0F3} 
\multicolumn{2}{|c|}{\multirow{-3}{*}{\cellcolor[HTML]{00C0F3}\textbf{\begin{tabular}[c]{@{}c@{}}Objetivo 2.5: Divulgar la producción de conocimientos, para que \\ pueda ser comprendida y usada\end{tabular}}}} \\ \hline
\rowcolor[HTML]{8ED8F8} 
\multicolumn{1}{|c|}{\cellcolor[HTML]{8ED8F8}\textbf{Acciones}} & \multicolumn{1}{c|}{\cellcolor[HTML]{8ED8F8}\textbf{\begin{tabular}[c]{@{}c@{}}Estrategias metodológicas \\ y tecnológicas.\end{tabular}}} \\ \hline
\multicolumn{1}{|l|}{Elaborar documentos formales} & \begin{tabular}[c]{@{}l@{}}Uso de software especializado para\\ redactar artículos cienctíficos.\end{tabular} \\ \hline
\multicolumn{1}{|l|}{\begin{tabular}[c]{@{}l@{}}Socializar los conocimientos mediante la \\ interacción con otros profesionales.\end{tabular}} & \begin{tabular}[c]{@{}l@{}}Asistencia a conferencias, congresos,\\ seminarios, etc.\end{tabular} \\ \hline
\multicolumn{1}{|l|}{\begin{tabular}[c]{@{}l@{}}Comunicar los resultados mediante un \\ lenguaje interdisciplinario.\end{tabular}} & Realizar informes y presentaciones. \\ \hline
\end{tabular}
\end{table}

\begin{center}
    \textbf{Objeto 3: Propiedades de los números.}
\end{center}

\begin{table}[H]
\begin{tabular}{|ll|}
\hline
\rowcolor[HTML]{00C0F3} 
\multicolumn{2}{|c|}{\cellcolor[HTML]{00C0F3}} \\
\rowcolor[HTML]{00C0F3} 
\multicolumn{2}{|c|}{\cellcolor[HTML]{00C0F3}} \\
\rowcolor[HTML]{00C0F3} 
\multicolumn{2}{|c|}{\multirow{-3}{*}{\cellcolor[HTML]{00C0F3}\textbf{\begin{tabular}[c]{@{}c@{}}Objetivo 3.1: Demostrar nuevos resultados para contribuir en el entorno\\ \\ cient'ifico y social.\end{tabular}}}} \\ \hline
\rowcolor[HTML]{8ED8F8} 
\multicolumn{1}{|c|}{\cellcolor[HTML]{8ED8F8}\textbf{Acciones}} & \multicolumn{1}{c|}{\cellcolor[HTML]{8ED8F8}\textbf{\begin{tabular}[c]{@{}c@{}}Estrategias metodológicas \\ \\ y tecnológicas\end{tabular}}} \\ \hline
\multicolumn{1}{|l|}{Conocer las definiciones y los principios básicos.} &  \\
\multicolumn{1}{|l|}{Generalizar resultados, nociones y definiciones.} & Exploración bibliográfica y web \\
\multicolumn{1}{|l|}{Relacionar distintas áreas de la Teoría de Números} & Resolución de ejercicios. \\
\multicolumn{1}{|l|}{Efectuar cálculos concretos.} & Trabajo en equipo. \\ \hline
\end{tabular}
\end{table}

\begin{table}[H]
\begin{tabular}{|ll|}
\hline
\rowcolor[HTML]{00C0F3} 
\multicolumn{2}{|c|}{\cellcolor[HTML]{00C0F3}} \\
\rowcolor[HTML]{00C0F3} 
\multicolumn{2}{|c|}{\cellcolor[HTML]{00C0F3}} \\
\rowcolor[HTML]{00C0F3} 
\multicolumn{2}{|c|}{\multirow{-3}{*}{\cellcolor[HTML]{00C0F3}\textbf{\begin{tabular}[c]{@{}c@{}}Objetivo 3.2: Conocer las propiedades de curvas, superficies y variedades para visualizar \\ las nociones geométricas básicas en espacios de dos, tres y más dimensiones.\\ .\end{tabular}}}} \\ \hline
\rowcolor[HTML]{8ED8F8} 
\multicolumn{1}{|c|}{\cellcolor[HTML]{8ED8F8}\textbf{Acciones}} & \multicolumn{1}{c|}{\cellcolor[HTML]{8ED8F8}\textbf{\begin{tabular}[c]{@{}c@{}}Estrategias metodológicas \\ \\ y tecnológicas\end{tabular}}} \\ \hline
\multicolumn{1}{|l|}{Elaborar algoritmos.} &  \\
\multicolumn{1}{|l|}{Analizar datos obtenidos.} & CAS \\
\multicolumn{1}{|l|}{\begin{tabular}[c]{@{}l@{}}Identificar patrones, tendencias y \\ comportamientos.\end{tabular}} & Lenguajes de programación. \\
\multicolumn{1}{|l|}{Implementar los algoritmos.} & \begin{tabular}[c]{@{}l@{}}Sistemas de cómputo especializados\\ en Teoría de Números.\end{tabular} \\ \hline
\rowcolor[HTML]{00C0F3} 
\multicolumn{2}{|c|}{\cellcolor[HTML]{00C0F3}} \\
\rowcolor[HTML]{00C0F3} 
\multicolumn{2}{|c|}{\cellcolor[HTML]{00C0F3}} \\
\rowcolor[HTML]{00C0F3} 
\multicolumn{2}{|c|}{\multirow{-3}{*}{\cellcolor[HTML]{00C0F3}\textbf{\begin{tabular}[c]{@{}c@{}}Objetivo 3.3: Formular conjeturas con base en observaciones experimentales\\ o teóricas sobre resultados plausibles para producir nuevo conocimiento.\end{tabular}}}} \\ \hline
\rowcolor[HTML]{8ED8F8} 
\multicolumn{1}{|c|}{\cellcolor[HTML]{8ED8F8}\textbf{Acciones}} & \multicolumn{1}{c|}{\cellcolor[HTML]{8ED8F8}\textbf{\begin{tabular}[c]{@{}c@{}}Estrategias metodológicas \\ y tecnológicas\end{tabular}}} \\ \hline
\multicolumn{1}{|l|}{Plantear problemas de estudio.} &  \\
\multicolumn{1}{|l|}{Generalizar resultados previamente encontrados.} & Exploración bibliográfica. \\
\multicolumn{1}{|l|}{Inducir resultados con base en la experimentación.} & \begin{tabular}[c]{@{}l@{}}Interacción con expertos de otras\\ áreas.\end{tabular} \\
\multicolumn{1}{|l|}{\begin{tabular}[c]{@{}l@{}}Ordenar resultados relacionados por medio de \\ puntos de vista unificadores.\end{tabular}} & Seminarios, charlas. \\ \hline
\rowcolor[HTML]{00C0F3} 
\multicolumn{2}{|c|}{\cellcolor[HTML]{00C0F3}} \\
\rowcolor[HTML]{00C0F3} 
\multicolumn{2}{|c|}{\cellcolor[HTML]{00C0F3}} \\
\rowcolor[HTML]{00C0F3} 
\multicolumn{2}{|c|}{\multirow{-3}{*}{\cellcolor[HTML]{00C0F3}\textbf{\begin{tabular}[c]{@{}c@{}}Objetivo 3.4: Demostrar nuevos resultados para contribuir en el\\ entorno científico y social.\end{tabular}}}} \\ \hline
\rowcolor[HTML]{8ED8F8} 
\multicolumn{1}{|c|}{\cellcolor[HTML]{8ED8F8}\textbf{Acciones}} & \multicolumn{1}{c|}{\cellcolor[HTML]{8ED8F8}\textbf{\begin{tabular}[c]{@{}c@{}}Estrategias metodológicas\\ \\  y tecnológicas\end{tabular}}} \\ \hline
\multicolumn{1}{|l|}{\begin{tabular}[c]{@{}l@{}}Estudiar modelos geométricos de problemas \\ reales abiertos.\end{tabular}} & \begin{tabular}[c]{@{}l@{}}Uso de software especializado como\\ CAS para ayudar en demostraciones \\ que involucran cálculos complicados.\end{tabular} \\
\multicolumn{1}{|l|}{\begin{tabular}[c]{@{}l@{}}Estudiar y aplicar las técnicas previas relacionadas\\ a los objetos estudiados.\end{tabular}} & \begin{tabular}[c]{@{}l@{}}Aplicación de técnicas usuales en el\\ estudio de los objetos particulares.\end{tabular} \\
\multicolumn{1}{|l|}{Crear nuevas teorías, conceptos y técnicas.} & \begin{tabular}[c]{@{}l@{}}Búsqueda de nuevos métodos para\\ solventar problemas nuevos.\end{tabular} \\
\multicolumn{1}{|l|}{Verificar y validar los teoremas obtenidos.} & \begin{tabular}[c]{@{}l@{}}Utilización de avances recientes para\\ demostrar nuevos teoremas.\end{tabular} \\
\multicolumn{1}{|l|}{} & \begin{tabular}[c]{@{}l@{}}Uso de software para verificar y \\ validar resultados.\end{tabular} \\ \hline
\rowcolor[HTML]{00C0F3} 
\multicolumn{2}{|c|}{\cellcolor[HTML]{00C0F3}} \\
\rowcolor[HTML]{00C0F3} 
\multicolumn{2}{|c|}{\cellcolor[HTML]{00C0F3}} \\
\rowcolor[HTML]{00C0F3} 
\multicolumn{2}{|c|}{\multirow{-3}{*}{\cellcolor[HTML]{00C0F3}\textbf{\begin{tabular}[c]{@{}c@{}}Objetivo 3.5: Divulgar la producción de conocimientos, para\\ \\ que pueda ser comprendida y usada.\end{tabular}}}} \\ \hline
\rowcolor[HTML]{8ED8F8} 
\multicolumn{1}{|c|}{\cellcolor[HTML]{8ED8F8}\textbf{Acciones}} & \multicolumn{1}{c|}{\cellcolor[HTML]{8ED8F8}\textbf{\begin{tabular}[c]{@{}c@{}}Estrategias metodológicas \\ y tecnológicas.\end{tabular}}} \\ \hline
\multicolumn{1}{|l|}{Elaborar documentos formales} & \begin{tabular}[c]{@{}l@{}}Uso de software especializado para\\ redactar artículos científicos.\end{tabular} \\
\multicolumn{1}{|l|}{\begin{tabular}[c]{@{}l@{}}Socializar los conocimientos mediante la \\ interacción con otros profesionales.\end{tabular}} & \begin{tabular}[c]{@{}l@{}}Asistencia a conferencias, congresos,\\ seminarios, etc.\end{tabular} \\
\multicolumn{1}{|l|}{\begin{tabular}[c]{@{}l@{}}Comunicar los resultados mediante un \\ lenguaje interdisciplinario.\end{tabular}} & Realizar informes y presentaciones. \\ \hline
\end{tabular}
\end{table}

\begin{center}
    \textbf{Objeto 4: Espacios Topológicos.}
\end{center}

\begin{table}[H]
\begin{tabular}{|cl|}
\hline
\rowcolor[HTML]{00C0F3} 
\multicolumn{2}{|c|}{\cellcolor[HTML]{00C0F3}} \\
\rowcolor[HTML]{00C0F3} 
\multicolumn{2}{|c|}{\cellcolor[HTML]{00C0F3}} \\
\rowcolor[HTML]{00C0F3} 
\multicolumn{2}{|c|}{\multirow{-3}{*}{\cellcolor[HTML]{00C0F3}\textbf{\begin{tabular}[c]{@{}c@{}}Objetivo 4.1: Entender qué es una Topología para que el estudiantado comprenda los \\ alcances de esta área de estudio y sus relaciones con otras áreas de la matemática\\ y la ciencia en general.\end{tabular}}}} \\ \hline
\rowcolor[HTML]{8ED8F8} 
\multicolumn{1}{|c|}{\cellcolor[HTML]{8ED8F8}\textbf{Acciones}} & \multicolumn{1}{c|}{\cellcolor[HTML]{8ED8F8}\textbf{\begin{tabular}[c]{@{}c@{}}Estrategias metodológicas \\ \\ y tecnológicas\end{tabular}}} \\ \hline
\multicolumn{1}{|l|}{Estudiar ejemplos de topología.} & Exploración bibliográfica y web. \\
\multicolumn{1}{|l|}{Estudiar propiedades de topología.} & Resolución de ejercicios. \\
\multicolumn{1}{|l|}{Clasificar los ejemplos de topologías.} & Trabajo en equipo. \\ \hline
\rowcolor[HTML]{00C0F3} 
\multicolumn{2}{|c|}{\cellcolor[HTML]{00C0F3}} \\
\rowcolor[HTML]{00C0F3} 
\multicolumn{2}{|c|}{\cellcolor[HTML]{00C0F3}} \\
\rowcolor[HTML]{00C0F3} 
\multicolumn{2}{|c|}{\multirow{-3}{*}{\cellcolor[HTML]{00C0F3}\textbf{\begin{tabular}[c]{@{}c@{}}Objetivo 4.2: Manipular las topologías para generar nuevos espacios topológicos\\ \\ y estudiar las relaciones entre nuevos y conocidos espacios topológicos.\end{tabular}}}} \\ \hline
\rowcolor[HTML]{8ED8F8} 
\multicolumn{1}{|c|}{\cellcolor[HTML]{8ED8F8}\textbf{Acciones}} & \multicolumn{1}{c|}{\cellcolor[HTML]{8ED8F8}\textbf{\begin{tabular}[c]{@{}c@{}}Estrategias metodológicas \\ \\ y tecnológicas\end{tabular}}} \\ \hline
\multicolumn{1}{|l|}{} & Exploración bibliográfica y web. \\
\multicolumn{1}{|l|}{Construir nuevas topologías a partir de otras.} & Resolución de ejercicios. \\
\multicolumn{1}{|l|}{Estudiar invariantes topológicos.} & Trabajo en equipo. \\ \hline
\rowcolor[HTML]{00C0F3} 
\multicolumn{2}{|c|}{\cellcolor[HTML]{00C0F3}} \\
\rowcolor[HTML]{00C0F3} 
\multicolumn{2}{|c|}{\cellcolor[HTML]{00C0F3}} \\
\rowcolor[HTML]{00C0F3} 
\multicolumn{2}{|c|}{\multirow{-3}{*}{\cellcolor[HTML]{00C0F3}\textbf{\begin{tabular}[c]{@{}c@{}}Objetivo 4.3: Incorporar otras áreas de la matemática en la topología para lograr\\ mejores herramientas para relacionar y distinguir los diferentes \\ espacios topológicos.\end{tabular}}}} \\ \hline
\rowcolor[HTML]{8ED8F8} 
\multicolumn{1}{|c|}{\cellcolor[HTML]{8ED8F8}\textbf{Acciones}} & \multicolumn{1}{c|}{\cellcolor[HTML]{8ED8F8}\textbf{\begin{tabular}[c]{@{}c@{}}Estrategias metodológicas \\ y tecnológicas\end{tabular}}} \\ \hline
\multicolumn{1}{|l|}{\begin{tabular}[c]{@{}l@{}}Incorporar el análisis como herramienta para la \\ topología (Topología métrica).\end{tabular}} & Exploración bibliográfica. \\
\multicolumn{1}{|l|}{\begin{tabular}[c]{@{}l@{}}Incorporar el Álgebra como herramienta para la\\ Topología (Topología diferencial).\end{tabular}} & \begin{tabular}[c]{@{}l@{}}Interacción con expertos de otras \\ áreas.\end{tabular} \\
\multicolumn{1}{|l|}{\begin{tabular}[c]{@{}l@{}}Incorporar variedades diferenciables como \\ herramientas la Topología (Topología \\ diferencial).\end{tabular}} & Seminarios, charlas. \\ \hline
\rowcolor[HTML]{00C0F3} 
\multicolumn{2}{|c|}{\cellcolor[HTML]{00C0F3}} \\
\rowcolor[HTML]{00C0F3} 
\multicolumn{2}{|c|}{\cellcolor[HTML]{00C0F3}} \\
\rowcolor[HTML]{00C0F3} 
\multicolumn{2}{|c|}{\multirow{-3}{*}{\cellcolor[HTML]{00C0F3}\textbf{\begin{tabular}[c]{@{}c@{}}Objetivo 4.4: Diferenciar espacios topológcos para clasificarlos\\ \\ y estudiar sus propiedades, similitudes y diferencias.\end{tabular}}}} \\ \hline
\rowcolor[HTML]{8ED8F8} 
\multicolumn{1}{|c|}{\cellcolor[HTML]{8ED8F8}\textbf{Acciones}} & \multicolumn{1}{c|}{\cellcolor[HTML]{8ED8F8}\textbf{\begin{tabular}[c]{@{}c@{}}Estrategias metodológicas\\ \\  y tecnológicas\end{tabular}}} \\ \hline
\multicolumn{1}{|l|}{} & Exploración bibliográfica. \\
\multicolumn{1}{|l|}{\begin{tabular}[c]{@{}l@{}}Estudiar propiedades de los espacios \\ topológicos.\end{tabular}} & \begin{tabular}[c]{@{}l@{}}Interacción con expertos de otras \\ áreas.\end{tabular} \\
\multicolumn{1}{|l|}{Estudiar invariantes topológicas.} & Seminarios, charlas. \\ \hline
\end{tabular}
\end{table}

\begin{table}[H]
\begin{tabular}{|cl|}
\hline
\rowcolor[HTML]{00C0F3} 
\multicolumn{2}{|c|}{\cellcolor[HTML]{00C0F3}} \\
\rowcolor[HTML]{00C0F3} 
\multicolumn{2}{|c|}{\cellcolor[HTML]{00C0F3}} \\
\rowcolor[HTML]{00C0F3} 
\multicolumn{2}{|c|}{\multirow{-3}{*}{\cellcolor[HTML]{00C0F3}\textbf{\begin{tabular}[c]{@{}c@{}}Objetivo 4.5: Aplicar la topología en otras áreas de la Matemática para generar\\ \\ nuevos resultados tanto en esta misma área como en otras áreas de \\ la matemática y la ciencia.\end{tabular}}}} \\ \hline
\rowcolor[HTML]{8ED8F8} 
\multicolumn{1}{|c|}{\cellcolor[HTML]{8ED8F8}\textbf{Acciones}} & \multicolumn{1}{c|}{\cellcolor[HTML]{8ED8F8}\textbf{\begin{tabular}[c]{@{}c@{}}Estrategias metodológicas \\ y tecnológicas.\end{tabular}}} \\ \hline
\multicolumn{1}{|l|}{} & Exploración bibliográfica. \\
\multicolumn{1}{|l|}{\begin{tabular}[c]{@{}l@{}}Estudiar teoremas de topología que se pueden \\ utilizar en otras áreas de la Matemática.\end{tabular}} & \begin{tabular}[c]{@{}l@{}}Interacción con expertos de otras\\ áreas.\end{tabular} \\
\multicolumn{1}{|l|}{\begin{tabular}[c]{@{}l@{}}Generar resultados en otras áreas de la \\ Matemática utilizando Topología.\end{tabular}} & Seminarios, charlas.. \\ \hline
\rowcolor[HTML]{00C0F3} 
\multicolumn{2}{|c|}{\cellcolor[HTML]{00C0F3}} \\
\rowcolor[HTML]{00C0F3} 
\multicolumn{2}{|c|}{\cellcolor[HTML]{00C0F3}} \\
\rowcolor[HTML]{00C0F3} 
\multicolumn{2}{|c|}{\multirow{-3}{*}{\cellcolor[HTML]{00C0F3}\textbf{\begin{tabular}[c]{@{}c@{}}Objetivo 4.6: Divulgar la producción de conocimientos, para \\ que pueda ser comprendida y usada.\end{tabular}}}} \\ \hline
\rowcolor[HTML]{8ED8F8} 
\multicolumn{1}{|c|}{\cellcolor[HTML]{8ED8F8}\textbf{Acciones}} & \multicolumn{1}{c|}{\cellcolor[HTML]{8ED8F8}\textbf{\begin{tabular}[c]{@{}c@{}}Estrategias metodológicas \\ y tecnológicas.\end{tabular}}} \\ \hline
\multicolumn{1}{|l|}{Elaborar documentos formales.} & \begin{tabular}[c]{@{}l@{}}Uso de software especializado para \\ redactar artículos científicos.\end{tabular} \\
\multicolumn{1}{|l|}{\begin{tabular}[c]{@{}l@{}}Socializar los conocimientos mediante la \\ interacción con otros profesionales.\end{tabular}} & \begin{tabular}[c]{@{}l@{}}Asistencia a conferencias, congresos, \\ seminarios, etc\end{tabular} \\
\multicolumn{1}{|l|}{\begin{tabular}[c]{@{}l@{}}Comunicar los resultados mediante un lenguaje \\ interdisciplinario.\end{tabular}} & Realizar informes y presentaciones. \\ \hline
\end{tabular}
\end{table}

\begin{center}
    \textbf{Objeto 5: Fenómenos, objetos o experimentos de carácter aleatorio}
\end{center}

\begin{table}[H]
\begin{tabular}{|cl|}
\hline
\rowcolor[HTML]{00C0F3} 
\multicolumn{2}{|c|}{\cellcolor[HTML]{00C0F3}} \\
\rowcolor[HTML]{00C0F3} 
\multicolumn{2}{|c|}{\cellcolor[HTML]{00C0F3}} \\
\rowcolor[HTML]{00C0F3} 
\multicolumn{2}{|c|}{\multirow{-3}{*}{\cellcolor[HTML]{00C0F3}\textbf{\begin{tabular}[c]{@{}c@{}}Objetivo 5.1: Comprender los objetos de naturaleza aleatoria para \\ realizar planteamientos matemáticos adecuados.\end{tabular}}}} \\ \hline
\rowcolor[HTML]{8ED8F8} 
\multicolumn{1}{|c|}{\cellcolor[HTML]{8ED8F8}\textbf{Acciones}} & \multicolumn{1}{c|}{\cellcolor[HTML]{8ED8F8}\textbf{\begin{tabular}[c]{@{}c@{}}Estrategias metodológicas \\ \\ y tecnológicas\end{tabular}}} \\ \hline
\multicolumn{1}{|l|}{Definir los objetos aleatorios.} & Investigar bibliografía y web \\ \hline
\multicolumn{1}{|l|}{\begin{tabular}[c]{@{}l@{}}Definir objetos matemáticos que resumen \\ información sobre fenómenos aleatorios\\ (medida de probabilidad, momentos, función\\ características, etc.).\end{tabular}} & Realización de prácticas. \\ \hline
\multicolumn{1}{|l|}{\begin{tabular}[c]{@{}l@{}}Estudiar la interacción entre fenómenos \\ aleatorios (independencias, correlación,etc.).\end{tabular}} & Exploración de datos. \\ \hline
\end{tabular}
\end{table}

\begin{table}[H]
\begin{tabular}{|cl|}
\hline
\rowcolor[HTML]{00C0F3} 
\multicolumn{2}{|c|}{\cellcolor[HTML]{00C0F3}} \\
\rowcolor[HTML]{00C0F3} 
\multicolumn{2}{|c|}{\cellcolor[HTML]{00C0F3}} \\
\rowcolor[HTML]{00C0F3} 
\multicolumn{2}{|c|}{\multirow{-3}{*}{\cellcolor[HTML]{00C0F3}\textbf{\begin{tabular}[c]{@{}c@{}}Objetivo 5.2: Estudiar los procesos estocásticos, con el fin de plantear \\ \\ adecuadamente los problemas .\end{tabular}}}} \\ \hline
\rowcolor[HTML]{8ED8F8} 
\multicolumn{1}{|c|}{\cellcolor[HTML]{8ED8F8}\textbf{Acciones}} & \multicolumn{1}{c|}{\cellcolor[HTML]{8ED8F8}\textbf{\begin{tabular}[c]{@{}c@{}}Estrategias metodológicas \\ \\ y tecnológicas\end{tabular}}} \\ \hline
\multicolumn{1}{|l|}{\begin{tabular}[c]{@{}l@{}}Distinguir los diferentes tipos de procesos\\ estocásticos.\end{tabular}} & \begin{tabular}[c]{@{}l@{}}Interacción con expertos (conferencias,\\ seminarios, charlas de otras disciplinas\\ y áreas de la matemática).\end{tabular} \\ \hline
\multicolumn{1}{|l|}{\begin{tabular}[c]{@{}l@{}}Caracterizar las propiedades de los diferentes \\ tipos de procesos estocásticos (convergencia,\\ regularidad de trayectorias, teoremas clásicos).\end{tabular}} & Experimentación con los datos. \\ \hline
\multicolumn{1}{|l|}{\begin{tabular}[c]{@{}l@{}}Construir nuevos procesos a apartit de procesos\\ existentes (integración estocástica, SPDE's, \\ SDE's, transformadas, etc.)\end{tabular}} & Consultas bibliográficas. \\ \hline
\rowcolor[HTML]{00C0F3} 
\multicolumn{2}{|c|}{\cellcolor[HTML]{00C0F3}} \\
\rowcolor[HTML]{00C0F3} 
\multicolumn{2}{|c|}{\cellcolor[HTML]{00C0F3}} \\
\rowcolor[HTML]{00C0F3} 
\multicolumn{2}{|c|}{\multirow{-3}{*}{\cellcolor[HTML]{00C0F3}\textbf{\begin{tabular}[c]{@{}c@{}}Objetivo 5.3: Modelar fenómenos aleatorios usando procesos estocásticos, para\\ posibilitar la etapa posterior de resolución del problema.\end{tabular}}}} \\ \hline
\rowcolor[HTML]{8ED8F8} 
\multicolumn{1}{|c|}{\cellcolor[HTML]{8ED8F8}\textbf{Acciones}} & \multicolumn{1}{c|}{\cellcolor[HTML]{8ED8F8}\textbf{\begin{tabular}[c]{@{}c@{}}Estrategias metodológicas \\ y tecnológicas\end{tabular}}} \\ \hline
\multicolumn{1}{|l|}{Identificar y decidir el tipo de proceso adecuado.} & \begin{tabular}[c]{@{}l@{}}Interacción con expertos (conferencias,\\ seminarios, charlas de otras disciplinas\\ y áreas de la matemática.)\end{tabular} \\ \hline
\multicolumn{1}{|l|}{\begin{tabular}[c]{@{}l@{}}Estudiar las propiedades usando teoría clásica o\\ formulando nueva teoría.\end{tabular}} & Experimentación con los datos. \\ \hline
\multicolumn{1}{|l|}{\begin{tabular}[c]{@{}l@{}}Inferir y comunicar resultados sobre el fenómeno\\ modelado.\end{tabular}} & Consultas bibliográficas. \\ \hline
\rowcolor[HTML]{00C0F3} 
\multicolumn{2}{|c|}{\cellcolor[HTML]{00C0F3}} \\
\rowcolor[HTML]{00C0F3} 
\multicolumn{2}{|c|}{\cellcolor[HTML]{00C0F3}} \\
\rowcolor[HTML]{00C0F3} 
\multicolumn{2}{|c|}{\multirow{-3}{*}{\cellcolor[HTML]{00C0F3}\textbf{\begin{tabular}[c]{@{}c@{}}Objetivo 5.4: Divulgar la producción de conocimientos, para que \\ pueda ser comprendida y usada.\end{tabular}}}} \\ \hline
\rowcolor[HTML]{8ED8F8} 
\multicolumn{1}{|c|}{\cellcolor[HTML]{8ED8F8}\textbf{Acciones}} & \multicolumn{1}{c|}{\cellcolor[HTML]{8ED8F8}\textbf{\begin{tabular}[c]{@{}c@{}}Estrategias metodológicas\\ \\  y tecnológicas\end{tabular}}} \\ \hline
\multicolumn{1}{|l|}{Elaborar documentos formales} & \begin{tabular}[c]{@{}l@{}}Uso de software especializado para\\ redactar artículos científicos.\end{tabular} \\ \hline
\multicolumn{1}{|l|}{\begin{tabular}[c]{@{}l@{}}Socializar los conocimientos mediante la\\ interacción con profesionales.\end{tabular}} & \begin{tabular}[c]{@{}l@{}}Asistencia a conferencias, congresos,\\ seminarios, etc.\end{tabular} \\ \hline
\multicolumn{1}{|l|}{\begin{tabular}[c]{@{}l@{}}Comunicar los resultados mediante un lenguaje\\ interdisciplinario.\end{tabular}} & Realizar informes y presentaciones. \\ \hline
\end{tabular}
\end{table}

\vspace{40mm}

\begin{center}
    \textbf{Objeto 6: Relación entre el lenguaje formal y los objetos matemáticos}
\end{center}

\begin{table}[H]
\begin{tabular}{|cl|}
\hline
\rowcolor[HTML]{00C0F3} 
\multicolumn{2}{|c|}{\cellcolor[HTML]{00C0F3}} \\
\rowcolor[HTML]{00C0F3} 
\multicolumn{2}{|c|}{\cellcolor[HTML]{00C0F3}} \\
\rowcolor[HTML]{00C0F3} 
\multicolumn{2}{|c|}{\multirow{-3}{*}{\cellcolor[HTML]{00C0F3}\textbf{\begin{tabular}[c]{@{}c@{}}Objetivo 6.1: Entender sintaxis formales para establecer lenguajes de\\ estudio de la matemática y de la informática.\end{tabular}}}} \\ \hline
\rowcolor[HTML]{8ED8F8} 
\multicolumn{1}{|c|}{\cellcolor[HTML]{8ED8F8}\textbf{Acciones}} & \multicolumn{1}{c|}{\cellcolor[HTML]{8ED8F8}\textbf{\begin{tabular}[c]{@{}c@{}}Estrategias metodológicas \\ \\ y tecnológicas\end{tabular}}} \\ \hline
\multicolumn{1}{|l|}{Estudiar cálculo proposicional.} & Exploración bibliográfica y web. \\
\multicolumn{1}{|l|}{Estudiar lenguajes de primer orden.} & Resolución de ejercicios. \\
\multicolumn{1}{|l|}{\begin{tabular}[c]{@{}l@{}}Ampliar el estudio a otros tipos de sintaxis \\ (lenguajes de segundo orden, de las Máquinas\\ de Turing, autómatas y lenguajes infinitarios).\end{tabular}} & Trabajo en equipo. \\ \hline
\rowcolor[HTML]{00C0F3} 
\multicolumn{2}{|c|}{\cellcolor[HTML]{00C0F3}} \\
\rowcolor[HTML]{00C0F3} 
\multicolumn{2}{|c|}{\cellcolor[HTML]{00C0F3}} \\
\rowcolor[HTML]{00C0F3} 
\multicolumn{2}{|c|}{\multirow{-3}{*}{\cellcolor[HTML]{00C0F3}\textbf{\begin{tabular}[c]{@{}c@{}}Objetivo 6.2: Entender la semántica en relación con la sintaxis para establecer\\ los modelos asociados a lenguajes matemáticos e informáticos.\end{tabular}}}} \\ \hline
\rowcolor[HTML]{8ED8F8} 
\multicolumn{1}{|c|}{\cellcolor[HTML]{8ED8F8}\textbf{Acciones}} & \multicolumn{1}{c|}{\cellcolor[HTML]{8ED8F8}\textbf{\begin{tabular}[c]{@{}c@{}}Estrategias metodológicas \\ \\ y tecnológicas\end{tabular}}} \\ \hline
\multicolumn{1}{|l|}{\begin{tabular}[c]{@{}l@{}}Construir la semántica en relación con la \\ sintaxis.\end{tabular}} & Exploración bibliográfica y web. \\
\multicolumn{1}{|l|}{\begin{tabular}[c]{@{}l@{}}Estudiar los teoremas de completitud y \\ compacidad.\end{tabular}} & Resolución de ejercicios. \\
\multicolumn{1}{|l|}{\begin{tabular}[c]{@{}l@{}}Estudiar aplicaciones de completitud y \\ compacidad en otras áreas de la Matemática\\ (Estructuras algebraicas, Teoría de números,\\ Análisis no estándar).\end{tabular}} & Trabajo en equipo. \\ \hline
\rowcolor[HTML]{00C0F3} 
\multicolumn{2}{|c|}{\cellcolor[HTML]{00C0F3}} \\
\rowcolor[HTML]{00C0F3} 
\multicolumn{2}{|c|}{\cellcolor[HTML]{00C0F3}} \\
\rowcolor[HTML]{00C0F3} 
\multicolumn{2}{|c|}{\multirow{-3}{*}{\cellcolor[HTML]{00C0F3}\textbf{\begin{tabular}[c]{@{}c@{}}Objetivo 6.3: Modelar fenómenos aleatorios usando procesos estocásticos, para\\ posibilitar la etapa posterior de resolución del problema.\end{tabular}}}} \\ \hline
\rowcolor[HTML]{8ED8F8} 
\multicolumn{1}{|c|}{\cellcolor[HTML]{8ED8F8}\textbf{Acciones}} & \multicolumn{1}{c|}{\cellcolor[HTML]{8ED8F8}\textbf{\begin{tabular}[c]{@{}c@{}}Estrategias metodológicas \\ y tecnológicas\end{tabular}}} \\ \hline
\multicolumn{1}{|l|}{} & Exploración bibliográfica \\
\multicolumn{1}{|l|}{\begin{tabular}[c]{@{}l@{}}Formalizar la computabilidad (funciones \\ recursivas, máquinas de Turing, entre otras).\end{tabular}} & \begin{tabular}[c]{@{}l@{}}Interacción con expertos de otras \\ áreas.\end{tabular} \\
\multicolumn{1}{|l|}{\begin{tabular}[c]{@{}l@{}}Mostrar el Teorema de Incompletitud de Godel\\ y el Teorema de Indecibilidad de Turing\end{tabular}} & Seminarios, charlas. \\ \hline
\rowcolor[HTML]{00C0F3} 
\multicolumn{2}{|c|}{\cellcolor[HTML]{00C0F3}} \\
\rowcolor[HTML]{00C0F3} 
\multicolumn{2}{|c|}{\cellcolor[HTML]{00C0F3}} \\
\rowcolor[HTML]{00C0F3} 
\multicolumn{2}{|c|}{\multirow{-3}{*}{\cellcolor[HTML]{00C0F3}\textbf{\begin{tabular}[c]{@{}c@{}}Objetivo 6.4: Estudiar las bases de la Teoría de Conjuntos para\\ dar consistencia a la matemática.\end{tabular}}}} \\ \hline
\rowcolor[HTML]{8ED8F8} 
\multicolumn{1}{|c|}{\cellcolor[HTML]{8ED8F8}\textbf{Acciones}} & \multicolumn{1}{c|}{\cellcolor[HTML]{8ED8F8}\textbf{\begin{tabular}[c]{@{}c@{}}Estrategias metodológicas\\ \\  y tecnológicas\end{tabular}}} \\ \hline
\multicolumn{1}{|l|}{\begin{tabular}[c]{@{}l@{}}Axiomatizar la teoría de conjuntos \\ (Zermelo-Fraenkel).\end{tabular}} & Exploración bibliográfica y web. \\
\multicolumn{1}{|l|}{Estudiar ordinales, cardinales y su aritmética.} & Resolución de ejercicios. \\
\multicolumn{1}{|l|}{\begin{tabular}[c]{@{}l@{}}Entender los modelos de ZF y sus extensiones\\ (ZFC).\end{tabular}} & Trabajo en equipo. \\
\multicolumn{1}{|l|}{Estudiar consistencia e independencia.} &  \\ \hline
\end{tabular}
\end{table}


\begin{table}[H]
\begin{tabular}{|cl|}
\hline
\rowcolor[HTML]{00C0F3} 
\multicolumn{2}{|c|}{\cellcolor[HTML]{00C0F3}} \\
\rowcolor[HTML]{00C0F3} 
\multicolumn{2}{|c|}{\cellcolor[HTML]{00C0F3}} \\
\rowcolor[HTML]{00C0F3} 
\multicolumn{2}{|c|}{\multirow{-3}{*}{\cellcolor[HTML]{00C0F3}\textbf{\begin{tabular}[c]{@{}c@{}}Objetivo 6.5: Producir conocimientos en la Lógica Matemática y áreas \\ \\ afines con el fin de aplicarlos a la matemática y a la informática.\end{tabular}}}} \\ \hline
\rowcolor[HTML]{8ED8F8} 
\multicolumn{1}{|c|}{\cellcolor[HTML]{8ED8F8}\textbf{Acciones}} & \multicolumn{1}{c|}{\cellcolor[HTML]{8ED8F8}\textbf{\begin{tabular}[c]{@{}c@{}}Estrategias metodológicas \\ y tecnológicas.\end{tabular}}} \\ \hline
\multicolumn{1}{|l|}{\begin{tabular}[c]{@{}l@{}}Profundizar en una o varias áreas de la Lógica\\ Matemática (Teoría de la Demostración de \\ Modelos, Computabilidad, Teoría de Conjuntos).\end{tabular}} & Exploración bibliográfica y web. \\
\multicolumn{1}{|l|}{Buscar nuevas preguntas de investigación.} & Resolución de ejercicios. \\
\multicolumn{1}{|l|}{Establecer nuevos resultados.} & Trabajo en equipo. \\ \hline
\multicolumn{1}{|l|}{\begin{tabular}[c]{@{}l@{}}Establecer relaciones con otras áreas de la \\ Matemática.\end{tabular}} &  \\ \hline
\rowcolor[HTML]{00C0F3} 
\multicolumn{2}{|c|}{\cellcolor[HTML]{00C0F3}} \\
\rowcolor[HTML]{00C0F3} 
\multicolumn{2}{|c|}{\cellcolor[HTML]{00C0F3}} \\
\rowcolor[HTML]{00C0F3} 
\multicolumn{2}{|c|}{\multirow{-3}{*}{\cellcolor[HTML]{00C0F3}\textbf{\begin{tabular}[c]{@{}c@{}}Objetivo 6.6: Divulgar la producción de conocimientos, para que pueda ser \\ comprendida y usada.\end{tabular}}}} \\ \hline
\rowcolor[HTML]{8ED8F8} 
\multicolumn{1}{|c|}{\cellcolor[HTML]{8ED8F8}\textbf{Acciones}} & \multicolumn{1}{c|}{\cellcolor[HTML]{8ED8F8}\textbf{\begin{tabular}[c]{@{}c@{}}Estrategias metodológicas \\ y tecnológicas.\end{tabular}}} \\ \hline
\multicolumn{1}{|l|}{Elaborar documentos formales.} & \begin{tabular}[c]{@{}l@{}}Uso de software especializado para \\ redactar artículos científicos.\end{tabular} \\
\multicolumn{1}{|l|}{\begin{tabular}[c]{@{}l@{}}Socializar los conocimientos mediante la \\ interacción con otros profesionales.\end{tabular}} & \begin{tabular}[c]{@{}l@{}}Asistencia a conferencias, congresos, \\ seminarios, etc.\end{tabular} \\
\multicolumn{1}{|l|}{\begin{tabular}[c]{@{}l@{}}Comunicar los resultados mediante un lenguaje \\ interdisciplinario.\end{tabular}} & Realizar informes y presentaciones. \\ \hline
\end{tabular}
\end{table}

\begin{center}
    \textbf{Objeto 7: Funciones a partir de relaciones entre sus derivadas (problema diferencial).}
\end{center}


\begin{table}[H]
\begin{tabular}{|cl|}
\hline
\rowcolor[HTML]{00C0F3} 
\multicolumn{2}{|c|}{\cellcolor[HTML]{00C0F3}} \\
\rowcolor[HTML]{00C0F3} 
\multicolumn{2}{|c|}{\cellcolor[HTML]{00C0F3}} \\
\rowcolor[HTML]{00C0F3} 
\multicolumn{2}{|c|}{\multirow{-3}{*}{\cellcolor[HTML]{00C0F3}\textbf{\begin{tabular}[c]{@{}c@{}}Objetivo 7.1: Determinar la naturaleza de los conceptos y relaciones en el\\ problema diferencial, con el fin de realizar un planteamiento\\ adecuado del problema.\end{tabular}}}} \\ \hline
\rowcolor[HTML]{8ED8F8} 
\multicolumn{1}{|c|}{\cellcolor[HTML]{8ED8F8}\textbf{Acciones}} & \multicolumn{1}{c|}{\cellcolor[HTML]{8ED8F8}\textbf{\begin{tabular}[c]{@{}c@{}}Estrategias metodológicas \\ y tecnológicas\end{tabular}}} \\ \hline
\multicolumn{1}{|l|}{} & Investigar bibliografía y web. \\
\multicolumn{1}{|l|}{\begin{tabular}[c]{@{}l@{}}Determinar los espacios de funciones \\ apropiados para el  problema.\end{tabular}} & Realización de prácticas. \\
\multicolumn{1}{|l|}{\begin{tabular}[c]{@{}l@{}}Definir el concepto de diferenciación y de \\ solución en estos espacios.\end{tabular}} & Exploración de datos. \\ \hline
\end{tabular}
\end{table}


\begin{table}[H]
\begin{tabular}{|cl|}
\hline
\rowcolor[HTML]{00C0F3} 
\multicolumn{2}{|c|}{\cellcolor[HTML]{00C0F3}} \\
\rowcolor[HTML]{00C0F3} 
\multicolumn{2}{|c|}{\cellcolor[HTML]{00C0F3}} \\
\rowcolor[HTML]{00C0F3} 
\multicolumn{2}{|c|}{\multirow{-3}{*}{\cellcolor[HTML]{00C0F3}\textbf{\begin{tabular}[c]{@{}c@{}}Objetivo 7.2: Describir cuantitativa o cualitativamente las soluciones\\ \\ del problema diferencial, con el fin de describir\\ o aproximar la solución exacta.\end{tabular}}}} \\ \hline
\rowcolor[HTML]{8ED8F8} 
\multicolumn{1}{|c|}{\cellcolor[HTML]{8ED8F8}\textbf{Acciones}} & \multicolumn{1}{c|}{\cellcolor[HTML]{8ED8F8}\textbf{\begin{tabular}[c]{@{}c@{}}Estrategias metodológicas \\ y tecnológicas\end{tabular}}} \\ \hline
\multicolumn{1}{|l|}{\begin{tabular}[c]{@{}l@{}}Estudiar la existencia y unicidad de las \\ soluciones.\end{tabular}} &  \\
\multicolumn{1}{|l|}{Estudiar la regularidad de las soluciones.} &  \\
\multicolumn{1}{|l|}{\begin{tabular}[c]{@{}l@{}}Estudiar la relación de las soluciones con sus\\ condiciones iniciales y de frontera (principio\\ máximo, método de energía, entre otros).\end{tabular}} & Investigar bibliografía y web. \\
\multicolumn{1}{|l|}{\begin{tabular}[c]{@{}l@{}}Estudiar el comportamiento asintótico y \\ estructural.\end{tabular}} & Realización de prácticas. \\
\multicolumn{1}{|l|}{Estudiar singularidades de solución.} & Exploración de datos. \\
\multicolumn{1}{|l|}{\begin{tabular}[c]{@{}l@{}}Estudiar métodos numéricos para la resolución\\ del problema diferencial.\end{tabular}} &  \\ \hline
\rowcolor[HTML]{00C0F3} 
\multicolumn{2}{|c|}{\cellcolor[HTML]{00C0F3}} \\
\rowcolor[HTML]{00C0F3} 
\multicolumn{2}{|c|}{\cellcolor[HTML]{00C0F3}} \\
\rowcolor[HTML]{00C0F3} 
\multicolumn{2}{|c|}{\multirow{-3}{*}{\cellcolor[HTML]{00C0F3}\textbf{\begin{tabular}[c]{@{}c@{}}Objetivo 7.3: Comprender las posibilidades de modelación de los fenómenos\\ diferenciales, para modelar y resolver problemas mediante el \\ uso de las ecuaciones diferenciales.\end{tabular}}}} \\ \hline
\rowcolor[HTML]{8ED8F8} 
\multicolumn{1}{|c|}{\cellcolor[HTML]{8ED8F8}\textbf{Acciones}} & \multicolumn{1}{c|}{\cellcolor[HTML]{8ED8F8}\textbf{\begin{tabular}[c]{@{}c@{}}Estrategias metodológicas \\ y tecnológicas\end{tabular}}} \\ \hline
\multicolumn{1}{|l|}{\begin{tabular}[c]{@{}l@{}}Entender la relación entre un problema \\ diferencial y el fenómeno que modela.\end{tabular}} &  \\
\multicolumn{1}{|l|}{\begin{tabular}[c]{@{}l@{}}Estudiar los problemas diferenciales clásicos\\ (sistemas lineales, elípticas, parabólicas, \\ hiperbólicas).\end{tabular}} & \begin{tabular}[c]{@{}l@{}}Interacción con expertos de otras \\ disciplinas.\end{tabular} \\
\multicolumn{1}{|l|}{\begin{tabular}[c]{@{}l@{}}Estudiar los problemas diferenciales \\ (linearización, análisis de Fourier,  cálculo \\ variacional, entre otros).\end{tabular}} & Trabajo interdisciplinario. \\
\multicolumn{1}{|l|}{Implementar métodos numéricos de simulación.} & Exploración numérica. \\ \hline
\rowcolor[HTML]{00C0F3} 
\multicolumn{2}{|c|}{\cellcolor[HTML]{00C0F3}} \\
\rowcolor[HTML]{00C0F3} 
\multicolumn{2}{|c|}{\cellcolor[HTML]{00C0F3}} \\
\rowcolor[HTML]{00C0F3} 
\multicolumn{2}{|c|}{\multirow{-3}{*}{\cellcolor[HTML]{00C0F3}\textbf{\begin{tabular}[c]{@{}c@{}}Objetivo 7.4: Producir nuevo conocimiento, apra contribuir en el entorno\\ \\ científico y social.\end{tabular}}}} \\ \hline
\rowcolor[HTML]{8ED8F8} 
\multicolumn{1}{|c|}{\cellcolor[HTML]{8ED8F8}\textbf{Acciones}} & \multicolumn{1}{c|}{\cellcolor[HTML]{8ED8F8}\textbf{\begin{tabular}[c]{@{}c@{}}Estrategias metodológicas\\  y tecnológicas\end{tabular}}} \\ \hline
\multicolumn{1}{|l|}{\begin{tabular}[c]{@{}l@{}}Entender y profundizar los problemas \\ diferenciales existentes.\end{tabular}} & Investigar bibliografía y web. \\
\multicolumn{1}{|l|}{\begin{tabular}[c]{@{}l@{}}Ampliar las posibilidades de modelación de los\\ problemas diferenciales.\end{tabular}} & Realización de prácticas. \\
\multicolumn{1}{|l|}{\begin{tabular}[c]{@{}l@{}}Establecer interacciones entre el área de \\ Ecuaciones Diferenciales y otras áreas del\\ conocimiento.\end{tabular}} & Exploración de datos. \\
\multicolumn{1}{|l|}{Estudiar consistencia e independencia.} &  \\ \hline
\end{tabular}
\end{table}


\begin{table}[H]
\begin{tabular}{|cl|}
\hline
\rowcolor[HTML]{00C0F3} 
\multicolumn{2}{|c|}{\cellcolor[HTML]{00C0F3}} \\
\rowcolor[HTML]{00C0F3} 
\multicolumn{2}{|c|}{\cellcolor[HTML]{00C0F3}} \\
\rowcolor[HTML]{00C0F3} 
\multicolumn{2}{|c|}{\multirow{-3}{*}{\cellcolor[HTML]{00C0F3}\textbf{\begin{tabular}[c]{@{}c@{}}Objetivo 7.5: Divulgar la producción de conocimientos, para que \\ \\ pueda ser comprendida y usada.\end{tabular}}}} \\ \hline
\rowcolor[HTML]{8ED8F8} 
\multicolumn{1}{|c|}{\cellcolor[HTML]{8ED8F8}\textbf{Acciones}} & \multicolumn{1}{c|}{\cellcolor[HTML]{8ED8F8}\textbf{\begin{tabular}[c]{@{}c@{}}Estrategias metodológicas \\ y tecnológicas.\end{tabular}}} \\ \hline
\multicolumn{1}{|l|}{Elaborar documentos formales.} & \begin{tabular}[c]{@{}l@{}}Uso de software especializado para\\ redactar artículos científicos.\end{tabular} \\
\multicolumn{1}{|l|}{\begin{tabular}[c]{@{}l@{}}Socializar los conocimientos mediante la \\ interacción con otros profesionales.\end{tabular}} & \begin{tabular}[c]{@{}l@{}}Asistencia a conferencias, congresos, \\ seminarios, etc.\end{tabular} \\
\multicolumn{1}{|l|}{\begin{tabular}[c]{@{}l@{}}Comunicar los resultados mediante un lenguaje\\ interdisciplinario.\end{tabular}} & Realizar informes y presentaciones. \\ \hline
\end{tabular}
\end{table}

\begin{center}
\textbf{Objeto 8: Estructuras asociadas a los números reales
o los números complejos, que involucren espacios topológicos.}
\end{center}

\begin{table}[H]
\begin{tabular}{|ll|}
\hline
\rowcolor[HTML]{00C0F3} 
\multicolumn{2}{|c|}{\cellcolor[HTML]{00C0F3}} \\
\rowcolor[HTML]{00C0F3} 
\multicolumn{2}{|c|}{\cellcolor[HTML]{00C0F3}} \\
\rowcolor[HTML]{00C0F3} 
\multicolumn{2}{|c|}{\multirow{-3}{*}{\cellcolor[HTML]{00C0F3}\textbf{\begin{tabular}[c]{@{}c@{}}Objetivo 8.1: Comprender el análisis sobre los espacios vectoriales reales de \\ \\ dimensión finita, con el fin de aplicarlo a situaciones dentro \\ fuera de la Matemática.\end{tabular}}}} \\ \hline
\rowcolor[HTML]{8ED8F8} 
\multicolumn{1}{|c|}{\cellcolor[HTML]{8ED8F8}\textbf{Acciones}} & \multicolumn{1}{c|}{\cellcolor[HTML]{8ED8F8}\textbf{\begin{tabular}[c]{@{}c@{}}Estrategias metodológicas \\ y tecnológicas\end{tabular}}} \\ \hline
\multicolumn{1}{|l|}{\begin{tabular}[c]{@{}l@{}}Estudiar técnicas de Cálculo Diferencial e \\ Integral en una y varias variables.\end{tabular}} & Exploración bibliográfica y web. \\
\multicolumn{1}{|l|}{\begin{tabular}[c]{@{}l@{}}Estudiar las propiedades topológicas de \\ espacios euclídeos.\end{tabular}} & Resolución de ejercicios. \\
\multicolumn{1}{|l|}{\begin{tabular}[c]{@{}l@{}}Estudiar las funciones, entre espacios euclídeos,\\ desde el punto de vista de su continuidad y \\ diferenciabilidad.\end{tabular}} & Trabajo en equipo. \\
\multicolumn{1}{|l|}{\begin{tabular}[c]{@{}l@{}}Estudiar los resultados mediante un lenguaje \\ interdisciplinario.\end{tabular}} &  \\ \hline
\rowcolor[HTML]{00C0F3} 
\multicolumn{2}{|c|}{\cellcolor[HTML]{00C0F3}} \\
\rowcolor[HTML]{00C0F3} 
\multicolumn{2}{|c|}{\cellcolor[HTML]{00C0F3}} \\
\rowcolor[HTML]{00C0F3} 
\multicolumn{2}{|c|}{\multirow{-3}{*}{\cellcolor[HTML]{00C0F3}\textbf{\begin{tabular}[c]{@{}c@{}}Objetivo 8.2: Comprender el análisis sobre los números complejos,\\ con el fin de aplicarlo a situaciones dentro o fuera de la Matemática.\end{tabular}}}} \\ \hline
\rowcolor[HTML]{8ED8F8} 
\multicolumn{1}{|c|}{\cellcolor[HTML]{8ED8F8}\textbf{Acciones}} & \multicolumn{1}{c|}{\cellcolor[HTML]{8ED8F8}\textbf{\begin{tabular}[c]{@{}c@{}}Estrategias metodológicas \\ y tecnológicas\end{tabular}}} \\ \hline
\multicolumn{1}{|l|}{\begin{tabular}[c]{@{}l@{}}Estudiar las funciones de variable compleja \\ desde el punto de vista de su continuidad y \\ diferenciabilidad.\end{tabular}} &  \\
\multicolumn{1}{|l|}{\begin{tabular}[c]{@{}l@{}}Estudiar las técnicas y resultados acerca de las\\ integrales de línea sobre los números complejos.\end{tabular}} & Exploración bibliográfica. \\
\multicolumn{1}{|l|}{\begin{tabular}[c]{@{}l@{}}Estudiar las propiedades de aproximación de\\ funciones de variable compleja.\end{tabular}} & \begin{tabular}[c]{@{}l@{}}Interacción con expertos de otras \\ áreas.\end{tabular} \\
\multicolumn{1}{|l|}{Estudiar las propiedades de mapeos conformes.} &  \\ \hline
\end{tabular}
\end{table}


\begin{table}[H]
\begin{tabular}{|ll|}
\hline
\rowcolor[HTML]{00C0F3} 
\multicolumn{2}{|c|}{\cellcolor[HTML]{00C0F3}} \\
\rowcolor[HTML]{00C0F3} 
\multicolumn{2}{|c|}{\cellcolor[HTML]{00C0F3}} \\
\rowcolor[HTML]{00C0F3} 
\multicolumn{2}{|c|}{\multirow{-3}{*}{\cellcolor[HTML]{00C0F3}\textbf{\begin{tabular}[c]{@{}c@{}}Objetivo 8.3: Conocer los principales conceptos y resultados del Análisis\\ Funcional, para abordar los problemas que involucran espacios de \\ dimensión infinita en su solución.\end{tabular}}}} \\ \hline
\rowcolor[HTML]{8ED8F8} 
\multicolumn{1}{|c|}{\cellcolor[HTML]{8ED8F8}\textbf{Acciones}} & \multicolumn{1}{c|}{\cellcolor[HTML]{8ED8F8}\textbf{\begin{tabular}[c]{@{}c@{}}Estrategias metodológicas \\ y tecnológicas\end{tabular}}} \\ \hline
\multicolumn{1}{|l|}{\begin{tabular}[c]{@{}l@{}}Estudiar las propiedades básicas de los\\ principales espacios de funciones \\ c,c, l,\end{tabular}} &  \\
\multicolumn{1}{|l|}{\begin{tabular}[c]{@{}l@{}}Estudiar las principales técnicas de aproximación \\ de espacios de funciones.\end{tabular}} &  \\
\multicolumn{1}{|l|}{\begin{tabular}[c]{@{}l@{}}Estudiar las propiedades algebraicas y topológicas\\ básicas de los espacios de dimensión infinita\\ (Banach, métricos, Fréchet, Hilbert).\end{tabular}} & Exploración bibliográfica. \\
\multicolumn{1}{|l|}{\begin{tabular}[c]{@{}l@{}}Estudiar las principales propiedades de las \\ transformaciones lineales continuas sobre\\ espacios vectoriales de dimensión infinita.\end{tabular}} & \begin{tabular}[c]{@{}l@{}}Interacción con expertos de otras\\ áreas.\end{tabular} \\
\multicolumn{1}{|l|}{\begin{tabular}[c]{@{}l@{}}Estudiar las propiedades básicas de los \\ operadores compactos.\end{tabular}} & Seminarios, charlas. \\
\multicolumn{1}{|l|}{\begin{tabular}[c]{@{}l@{}}Estudiar los rudimentos de la teoría espectral de \\ operadores.\end{tabular}} &  \\ \hline
\rowcolor[HTML]{00C0F3} 
\multicolumn{2}{|c|}{\cellcolor[HTML]{00C0F3}} \\
\rowcolor[HTML]{00C0F3} 
\multicolumn{2}{|c|}{\cellcolor[HTML]{00C0F3}} \\
\rowcolor[HTML]{00C0F3} 
\multicolumn{2}{|c|}{\multirow{-3}{*}{\cellcolor[HTML]{00C0F3}\textbf{\begin{tabular}[c]{@{}c@{}}Objetivo 8.4:  Investigar las propiedades de las estructuras del Análisis \\ \\ para generar nuevo conocimiento.\end{tabular}}}} \\ \hline
\rowcolor[HTML]{8ED8F8} 
\multicolumn{1}{|c|}{\cellcolor[HTML]{8ED8F8}\textbf{Acciones}} & \multicolumn{1}{c|}{\cellcolor[HTML]{8ED8F8}\textbf{\begin{tabular}[c]{@{}c@{}}Estrategias metodológicas\\  y tecnológicas\end{tabular}}} \\ \hline
\multicolumn{1}{|l|}{Revisar las propiedades conocidas.} &  \\
\multicolumn{1}{|l|}{Establecer conjeturas para nuevas propiedades.} & Exploración bibliográfica. \\
\multicolumn{1}{|l|}{\begin{tabular}[c]{@{}l@{}}Indagar los métodos y técnicas para resolver las\\ conjeturas.\end{tabular}} & \begin{tabular}[c]{@{}l@{}}Interacción con expertos de otras\\ áreas.\end{tabular} \\
\multicolumn{1}{|l|}{Resolver los problemas planteados.} & Seminarios, charlas. \\
\multicolumn{1}{|l|}{\begin{tabular}[c]{@{}l@{}}Identificar aplicaciones a otros problemas en \\ Análisis, otras áreas de la Matemática u otras\\ áreas del conocimiento.\end{tabular}} &  \\ \hline
\rowcolor[HTML]{00C0F3} 
\multicolumn{2}{|c|}{\cellcolor[HTML]{00C0F3}} \\
\rowcolor[HTML]{00C0F3} 
\multicolumn{2}{|c|}{\cellcolor[HTML]{00C0F3}} \\
\rowcolor[HTML]{00C0F3} 
\multicolumn{2}{|c|}{\multirow{-3}{*}{\cellcolor[HTML]{00C0F3}\textbf{\begin{tabular}[c]{@{}c@{}}Objetivo 8.5: Divulgar la producción de conocimientos, para que \\ pueda ser comprendida y usada.\end{tabular}}}} \\ \hline
\rowcolor[HTML]{8ED8F8} 
\multicolumn{1}{|c|}{\cellcolor[HTML]{8ED8F8}\textbf{Acciones}} & \multicolumn{1}{c|}{\cellcolor[HTML]{8ED8F8}\textbf{\begin{tabular}[c]{@{}c@{}}Estrategias metodológicas \\ y tecnológicas.\end{tabular}}} \\ \hline
\multicolumn{1}{|l|}{Elaborar documentos formales.} & \begin{tabular}[c]{@{}l@{}}Uso de software especializado para\\ redactar artículos científicos.\end{tabular} \\
\multicolumn{1}{|l|}{\begin{tabular}[c]{@{}l@{}}Socializar los conocimientos mediante la \\ interacción con otros profesionales.\end{tabular}} & \begin{tabular}[c]{@{}l@{}}Asistencia a conferencias, congresos, \\ seminarios, etc.\end{tabular} \\
\multicolumn{1}{|l|}{\begin{tabular}[c]{@{}l@{}}Comunicar los resultados mediante un lenguaje\\ interdisciplinario.\end{tabular}} & Realizar informes y presentaciones. \\ \hline
\end{tabular}
\end{table}

\vspace{16mm}

\begin{center}
    \textbf{Objeto 9: Algoritmos para aproximar soluciones de problemas matemáticos de diversas áreas.}
\end{center}

\begin{table}[H]
\begin{tabular}{|cl|}
\hline
\rowcolor[HTML]{00C0F3} 
\multicolumn{2}{|c|}{\cellcolor[HTML]{00C0F3}} \\
\rowcolor[HTML]{00C0F3} 
\multicolumn{2}{|c|}{\cellcolor[HTML]{00C0F3}} \\
\rowcolor[HTML]{00C0F3} 
\multicolumn{2}{|c|}{\multirow{-3}{*}{\cellcolor[HTML]{00C0F3}\textbf{\begin{tabular}[c]{@{}c@{}}Objetivo 9.1:  Comprender los algortimos clásicos de aproximación para aplicarlos \\ \\ en la solución de problemas.\end{tabular}}}} \\ \hline
\rowcolor[HTML]{8ED8F8} 
\multicolumn{1}{|c|}{\cellcolor[HTML]{8ED8F8}\textbf{Acciones}} & \multicolumn{1}{c|}{\cellcolor[HTML]{8ED8F8}\textbf{\begin{tabular}[c]{@{}c@{}}Estrategias metodológicas \\ y tecnológicas\end{tabular}}} \\ \hline
\multicolumn{1}{|l|}{} & Investigación bibliográfica. \\
\multicolumn{1}{|l|}{\begin{tabular}[c]{@{}l@{}}Estudiar características propias de algoritmos \\ existentes.\end{tabular}} & Uso de tecnología. \\
\multicolumn{1}{|l|}{Comparar los algoritmos.} & Resolución de ejercicios. \\
\multicolumn{1}{|l|}{\begin{tabular}[c]{@{}l@{}}Implementar algoritmos y comparar \\ rendimientos\end{tabular}} & Interacción con expertos. \\ \hline
\rowcolor[HTML]{00C0F3} 
\multicolumn{2}{|c|}{\cellcolor[HTML]{00C0F3}} \\
\rowcolor[HTML]{00C0F3} 
\multicolumn{2}{|c|}{\cellcolor[HTML]{00C0F3}} \\
\rowcolor[HTML]{00C0F3} 
\multicolumn{2}{|c|}{\multirow{-3}{*}{\cellcolor[HTML]{00C0F3}\textbf{\begin{tabular}[c]{@{}c@{}}Objetivo 9.2: Aplicar algoritmos clásicos para analizar fenómenos de\\ \\ diversas áreas.\end{tabular}}}} \\ \hline
\rowcolor[HTML]{8ED8F8} 
\multicolumn{1}{|c|}{\cellcolor[HTML]{8ED8F8}\textbf{Acciones}} & \multicolumn{1}{c|}{\cellcolor[HTML]{8ED8F8}\textbf{\begin{tabular}[c]{@{}c@{}}Estrategias metodológicas \\ y tecnológicas\end{tabular}}} \\ \hline
\multicolumn{1}{|l|}{} & Investigación bibliográfica. \\
\multicolumn{1}{|l|}{Identificar la naturaleza del fenómeno de estudio.} & Uso de tecnología. \\
\multicolumn{1}{|l|}{\begin{tabular}[c]{@{}l@{}}Implementar y comparar diferentes algoritmos que\\ resuelvan el problema.\end{tabular}} & Elaboración de reportes. \\
\multicolumn{1}{|l|}{Interpretar resultados obtenidos.} & Interacción con expertos. \\ \hline
\rowcolor[HTML]{00C0F3} 
\multicolumn{2}{|c|}{\cellcolor[HTML]{00C0F3}} \\
\rowcolor[HTML]{00C0F3} 
\multicolumn{2}{|c|}{\cellcolor[HTML]{00C0F3}} \\
\rowcolor[HTML]{00C0F3} 
\multicolumn{2}{|c|}{\multirow{-3}{*}{\cellcolor[HTML]{00C0F3}\textbf{\begin{tabular}[c]{@{}c@{}}Objetivo 9.3: Producir o mejorar algoritmos para aproximar soluciones\\ \\ de problemas numéricos.\end{tabular}}}} \\ \hline
\rowcolor[HTML]{8ED8F8} 
\multicolumn{1}{|c|}{\cellcolor[HTML]{8ED8F8}\textbf{Acciones}} & \multicolumn{1}{c|}{\cellcolor[HTML]{8ED8F8}\textbf{\begin{tabular}[c]{@{}c@{}}Estrategias metodológicas \\ y tecnológicas\end{tabular}}} \\ \hline
\multicolumn{1}{|l|}{Proponer algoritmos.} & Investigación bibliográfica. \\
\multicolumn{1}{|l|}{\begin{tabular}[c]{@{}l@{}}Verificar convergencia, estabilidad, planteamientos\\ y estimados de error.\end{tabular}} & Uso de tecnología. \\
\multicolumn{1}{|l|}{\begin{tabular}[c]{@{}l@{}}Implementar y comparar eficiencia (tiempos\\ de ejecución, complejidad, convergencia).\end{tabular}} & Interacción con expertos. \\ \hline
\rowcolor[HTML]{00C0F3} 
\multicolumn{2}{|c|}{\cellcolor[HTML]{00C0F3}} \\
\rowcolor[HTML]{00C0F3} 
\multicolumn{2}{|c|}{\cellcolor[HTML]{00C0F3}} \\
\rowcolor[HTML]{00C0F3} 
\multicolumn{2}{|c|}{\multirow{-3}{*}{\cellcolor[HTML]{00C0F3}\textbf{\begin{tabular}[c]{@{}c@{}}Objetivo 9.4: Divulgar la producción de conocimientos, para que \\ pueda ser comprendida y usada.\end{tabular}}}} \\ \hline
\rowcolor[HTML]{8ED8F8} 
\multicolumn{1}{|c|}{\cellcolor[HTML]{8ED8F8}\textbf{Acciones}} & \multicolumn{1}{c|}{\cellcolor[HTML]{8ED8F8}\textbf{\begin{tabular}[c]{@{}c@{}}Estrategias metodológicas \\ y tecnológicas.\end{tabular}}} \\ \hline
\multicolumn{1}{|l|}{Elaborar documentos formales.} & \begin{tabular}[c]{@{}l@{}}Uso de software especializado para\\ redactar artículos científicos.\end{tabular} \\
\multicolumn{1}{|l|}{\begin{tabular}[c]{@{}l@{}}Socializar los conocimientos mediante la \\ interacción con otros profesionales.\end{tabular}} & \begin{tabular}[c]{@{}l@{}}Asistencia a conferencias, congresos, \\ seminarios, etc.\end{tabular} \\
\multicolumn{1}{|l|}{\begin{tabular}[c]{@{}l@{}}Comunicar los resultados mediante un lenguaje\\ interdisciplinario.\end{tabular}} & Realizar informes y presentaciones. \\ \hline
\end{tabular}
\end{table}

\vspace{20mm}

\begin{center}
    \textbf{Objeto 10: Procesos de la vida real.}
\end{center}

\begin{table}[H]
\begin{tabular}{|cl|}
\hline
\rowcolor[HTML]{00C0F3} 
\multicolumn{2}{|c|}{\cellcolor[HTML]{00C0F3}} \\
\rowcolor[HTML]{00C0F3} 
\multicolumn{2}{|c|}{\cellcolor[HTML]{00C0F3}} \\
\rowcolor[HTML]{00C0F3} 
\multicolumn{2}{|c|}{\multirow{-3}{*}{\cellcolor[HTML]{00C0F3}\textbf{Objetivo 10.1: Indagar la especificidad de un proceso.}}} \\ \hline
\rowcolor[HTML]{8ED8F8} 
\multicolumn{1}{|c|}{\cellcolor[HTML]{8ED8F8}\textbf{Acciones}} & \multicolumn{1}{c|}{\cellcolor[HTML]{8ED8F8}\textbf{\begin{tabular}[c]{@{}c@{}}Estrategias metodológicas \\ y tecnológicas\end{tabular}}} \\ \hline
\multicolumn{1}{|l|}{Entender sobre el área de estudio.} &  \\
\multicolumn{1}{|l|}{Plantear problemas o preguntas.} & Exploración bibliográfica y web. \\
\multicolumn{1}{|l|}{Recolectar datos disponibles.} & Resolución de ejercicios. \\
\multicolumn{1}{|l|}{Identificar variables y parámetros.} & Trabajo en equipo. \\
\multicolumn{1}{|l|}{Identificar limitaciones y realizar suposiciones.} &  \\ \hline
\rowcolor[HTML]{00C0F3} 
\multicolumn{2}{|c|}{\cellcolor[HTML]{00C0F3}} \\
\rowcolor[HTML]{00C0F3} 
\multicolumn{2}{|c|}{\cellcolor[HTML]{00C0F3}} \\
\rowcolor[HTML]{00C0F3} 
\multicolumn{2}{|c|}{\multirow{-3}{*}{\cellcolor[HTML]{00C0F3}\textbf{Objetivo 10.2: Describir el proceso de la vida real con ecuaciones matemáticas.}}} \\ \hline
\rowcolor[HTML]{8ED8F8} 
\multicolumn{1}{|c|}{\cellcolor[HTML]{8ED8F8}\textbf{Acciones}} & \multicolumn{1}{c|}{\cellcolor[HTML]{8ED8F8}\textbf{\begin{tabular}[c]{@{}c@{}}Estrategias metodológicas \\ y tecnológicas\end{tabular}}} \\ \hline
\multicolumn{1}{|l|}{\begin{tabular}[c]{@{}l@{}}Determinar la herramienta matemática pertinente\\ con la cual se modelará.\end{tabular}} &  \\
\multicolumn{1}{|l|}{Formular el modelo con base en la herramienta.} & Exploración bibliográfica. \\
\multicolumn{1}{|l|}{Calibrar el modelo.} & \begin{tabular}[c]{@{}l@{}}Interacción con expertos de otras \\ áreas.\end{tabular} \\
\multicolumn{1}{|l|}{Reformular el modelo.} &  \\ \hline
\rowcolor[HTML]{00C0F3} 
\multicolumn{2}{|c|}{\cellcolor[HTML]{00C0F3}} \\
\rowcolor[HTML]{00C0F3} 
\multicolumn{2}{|c|}{\cellcolor[HTML]{00C0F3}} \\
\rowcolor[HTML]{00C0F3} 
\multicolumn{2}{|c|}{\multirow{-3}{*}{\cellcolor[HTML]{00C0F3}\textbf{Objetivo 10.3: Implementar el modelo planteado.}}} \\ \hline
\rowcolor[HTML]{8ED8F8} 
\multicolumn{1}{|c|}{\cellcolor[HTML]{8ED8F8}\textbf{Acciones}} & \multicolumn{1}{c|}{\cellcolor[HTML]{8ED8F8}\textbf{\begin{tabular}[c]{@{}c@{}}Estrategias metodológicas \\ y tecnológicas\end{tabular}}} \\ \hline
\multicolumn{1}{|l|}{Aplicar la teoría matemática al modelo.} & Exploración bibliográfica. \\
\multicolumn{1}{|l|}{Categorizar y organizar los resultados.} & \begin{tabular}[c]{@{}l@{}}Interacción con expertos de otras\\ áreas.\end{tabular} \\
\multicolumn{1}{|l|}{Formular recomendaciones.} & Seminarios, charlas. \\ \hline
\rowcolor[HTML]{00C0F3} 
\multicolumn{2}{|c|}{\cellcolor[HTML]{00C0F3}} \\
\rowcolor[HTML]{00C0F3} 
\multicolumn{2}{|c|}{\cellcolor[HTML]{00C0F3}} \\
\rowcolor[HTML]{00C0F3} 
\multicolumn{2}{|c|}{\multirow{-3}{*}{\cellcolor[HTML]{00C0F3}\textbf{Objetivo 10.4: Analizar los resultados del modelo planteado para el proceso.}}} \\ \hline
\rowcolor[HTML]{8ED8F8} 
\multicolumn{1}{|c|}{\cellcolor[HTML]{8ED8F8}\textbf{Acciones}} & \multicolumn{1}{c|}{\cellcolor[HTML]{8ED8F8}\textbf{\begin{tabular}[c]{@{}c@{}}Estrategias metodológicas\\  y tecnológicas\end{tabular}}} \\ \hline
\multicolumn{1}{|l|}{Interpretar los resultados con base en la pregunta.} & Software especializado. \\
\multicolumn{1}{|l|}{Determinar estrategias de abordaje al problema.} & Interacción con otras disciplinas. \\ \hline
\end{tabular}
\end{table}

\begin{table}[H]
\begin{tabular}{|cl|}
\hline
\rowcolor[HTML]{00C0F3} 
\multicolumn{2}{|c|}{\cellcolor[HTML]{00C0F3}} \\
\rowcolor[HTML]{00C0F3} 
\multicolumn{2}{|c|}{\cellcolor[HTML]{00C0F3}} \\
\rowcolor[HTML]{00C0F3} 
\multicolumn{2}{|c|}{\multirow{-3}{*}{\cellcolor[HTML]{00C0F3}\textbf{\begin{tabular}[c]{@{}c@{}}Objetivo 10.5: Divulgar la producción de conocimientos, para que \\ pueda ser comprendida y usada.\end{tabular}}}} \\ \hline
\rowcolor[HTML]{8ED8F8} 
\multicolumn{1}{|c|}{\cellcolor[HTML]{8ED8F8}\textbf{Acciones}} & \multicolumn{1}{c|}{\cellcolor[HTML]{8ED8F8}\textbf{\begin{tabular}[c]{@{}c@{}}Estrategias metodológicas \\ y tecnológicas.\end{tabular}}} \\ \hline
\multicolumn{1}{|l|}{Elaborar documentos formales.} & \begin{tabular}[c]{@{}l@{}}Uso de software especializado para\\ redactar artículos científicos.\end{tabular} \\
\multicolumn{1}{|l|}{\begin{tabular}[c]{@{}l@{}}Socializar los conocimientos mediante la \\ interacción con otros profesionales.\end{tabular}} & \begin{tabular}[c]{@{}l@{}}Asistencia a conferencias, congresos, \\ seminarios, etc.\end{tabular} \\
\multicolumn{1}{|l|}{\begin{tabular}[c]{@{}l@{}}Comunicar los resultados mediante un lenguaje\\ interdisciplinario.\end{tabular}} & Realizar informes y presentaciones. \\ \hline
\end{tabular}
\end{table}


\section{Perfil de salida}

El perfil de salida fue aprobado por la Asamblea de la Escuela de Matemática en Diciembre del 2020.

Este documento presenta una propuesta estructurada del perfil de egreso para la carrera de Bach. y Lic. en Matemática de la Universidad de Costa Rica, como parte de un proceso más amplio de reestructuración curricular. Dicho proceso surge del interés institucional por actualizar las carreras de Matemática y Ciencias Actuariales y se formaliza a través de una subcomisión de la Comisión de Docencia, con acompañamiento del Centro de Evaluación Académica (CEA). El documento constituye el primer paso para fundamentar la renovación del plan de estudios de la carrera, a partir del establecimiento de un perfil profesional coherente con las necesidades del siglo XXI.

La construcción del perfil implicó un trabajo colectivo y sistemático entre 2018 y 2020, que integró entrevistas a docentes, grupos focales, talleres, consultas por correo y análisis estructurales de la profesión matemática. Estos esfuerzos condujeron a la identificación de cinco áreas curriculares clave: habilidades operacionales, pensamiento abstracto, formalismo, pensamiento algorítmico y habilidades investigativas.

El documento se organiza en varios apartados. La Metodología detalla el proceso de trabajo, los actores involucrados y las estrategias empleadas para la recolección y análisis de información. Los Marcos referenciales abarcan un análisis histórico, profesional, disciplinar y pedagógico de la matemática como ciencia y como profesión. La Justificación argumenta la pertinencia de actualizar la formación matemática frente a los retos actuales de la sociedad y la ciencia, con respaldo en organismos internacionales como la UNESCO. Finalmente, se presentan una declaración de propósitos y el perfil de salida esperado de la persona graduada, perfil que incluye conocimientos, capacidades y actitudes fundamentales para su inserción profesional, académica o investigativa.

Este perfil busca dotar a la carrera de una identidad clara y de una base para el diseño curricular, articulando la tradición matemática con los desafíos contemporáneos en educación superior.

Las siguientes tablas contienen, de manera resumida, los aspectos fundamentales, incluyendo el código y el nivel, de cada saber del perfil de salida. Estos se dividen en los saberes conocer, los saberes hacer y los saberes ser.


\begin{longtable}{|l|l|l|}
\hline
\rowcolor[HTML]{00C0F3} 
\multicolumn{1}{|c|}{\cellcolor[HTML]{00C0F3}\textbf{Código}} & \multicolumn{1}{c|}{\cellcolor[HTML]{00C0F3}\textbf{Tipo del saber}} & \multicolumn{1}{c|}{\cellcolor[HTML]{00C0F3}\textbf{Descripción del saber}} \\ \hline
SC01 & Saber Conocer & \begin{tabular}[c]{@{}l@{}}Conocimiento avanzado en los fundamentos y\\  conceptos básicos de la matemática: operadores\\ lógicos, predicados, proposiciones, variables, \\ aritmética, la geometría euclídea.\end{tabular} \\ \hline
SC02 & Saber Conocer & \begin{tabular}[c]{@{}l@{}}Conocimiento avanzado en software avanzado\\ y aspectos computacionales de la matemática.\end{tabular} \\ \hline
SC03 & Saber Conocer & \begin{tabular}[c]{@{}l@{}}Conocimiento básico en lenguajes de programa-\\ ción que potencia su capacidad para programar\\ algoritmos matemáticos.\end{tabular} \\ \hline
SC04 & Saber Conocer & \begin{tabular}[c]{@{}l@{}}Conocimiento avanzado en técnicas de cálculo \\ en una y varias variables, y álgebra lineal.\end{tabular} \\ \hline
SC05 & Saber Conocer & \begin{tabular}[c]{@{}l@{}}Conocimiento intermedio en ecuaciones diferenciales,\\ teoría de la medida y probabilidad.\end{tabular} \\ \hline
SC06 & Saber Conocer & \begin{tabular}[c]{@{}l@{}}Conocimiento avanzado en el análisis de espacios\\ vectoriales reales de dimensión finita.\end{tabular} \\ \hline
SC07 & Saber Conocer & \begin{tabular}[c]{@{}l@{}}Conocimiento intermedio en los principales tipos\\ de Geometría (Euclídea, Proyectiva, Hiperbólica,\\ de curvas y superficies).\end{tabular} \\ \hline
SC08 & Saber Conocer & \begin{tabular}[c]{@{}l@{}}Conocimiento básico en los objetos de naturaleza\\ aleatoria (probabilidad, procesos estocásticos) y\\ sus aplicaciones en el entorno de la vida real.\end{tabular} \\ \hline
SC09 & Saber Conocer & \begin{tabular}[c]{@{}l@{}}Conocimiento básico en el análisis sobre los \\ números complejos y sus aplicaciones en el \\ entorno de la vida real.\end{tabular} \\ \hline
SC10 & Saber Conocer & \begin{tabular}[c]{@{}l@{}}Conocimiento intermedio en las estructuras al-\\ gebraicas(álgebra abstracta) y sus morfismos \\ junto con las aplicaciones a otras áreas del \\ conocimiento y del entorno de la vida real.\end{tabular} \\ \hline
SC11 & Saber Conocer & \begin{tabular}[c]{@{}l@{}}Conocimiento básico en las estructuras geomé-\\ tricas desde el punto de vista algebraico o dife-\\ rencial junto con sus aplicaciones a otras áreas\\ del conocimiento y el entorno de la vida real.\end{tabular} \\ \hline
SC12 & Saber Conocer & \begin{tabular}[c]{@{}l@{}}Conocimiento básico en la lógica matemática y \\ los fundamentos de la matemática, así como sus\\ aplicaciones a otras áreas del conocimiento y del \\ entorno de la vida real.\end{tabular} \\ \hline
SC13 & Saber Conocer & \begin{tabular}[c]{@{}l@{}}Conocimiento básico de la estructura de los sis-\\ temas numéricos (análisis numérico)  y sus apli-\\ caciones en el entorno y en otras disciplinas.\end{tabular} \\ \hline
SC14 & Saber Conocer & \begin{tabular}[c]{@{}l@{}}Conocimiento básico en los principales conceptos\\ y resultados del Análisis Funcional y sus aplica-\\ ciones a otras áreas del conocimiento.\end{tabular} \\ \hline
SC15 & Saber Conocer & \begin{tabular}[c]{@{}l@{}}Conocimiento intermedio en las estructuras topo-\\ lógicas desde el punto de vista general, geométrico\\ y algebraico, junto con las aplicaciones a otras \\ áreas de la matemática y de la vida real.\end{tabular} \\ \hline
SC16 & Saber Conocer & \begin{tabular}[c]{@{}l@{}}Conocimiento básico de los objetos de la Teoría\\ de Números y sus relaciones con otras áreas del \\ conocimiento.\end{tabular} \\ \hline
SC17 & Saber Conocer & \begin{tabular}[c]{@{}l@{}}Conocimiento básico en modelación matemática y\\ sus aplicaciones a la vida del entorno real.\end{tabular} \\ \hline
SH01 & Saber Hacer & \begin{tabular}[c]{@{}l@{}}Piensa operacionalmente en áreas de la matemá-\\ tica básica y avanzada para el trabajo propio y \\ con profesionales de otras disciplinas.\end{tabular} \\ \hline
SH02 & Saber Hacer & \begin{tabular}[c]{@{}l@{}}Resuelve problemas elementales de aritmética,\\ álgebra y cálculo.\end{tabular} \\ \hline
SH03 & Saber Hacer & \begin{tabular}[c]{@{}l@{}}Experimenta con nuevas tecnologías para validar \\ y conjeturar resultados.\end{tabular} \\ \hline
SH04 & Saber Hacer & \begin{tabular}[c]{@{}l@{}}Aplica la matemática a la búsqueda de soluciones\\ del entorno\end{tabular} \\ \hline
SH05 & Saber Hacer & \begin{tabular}[c]{@{}l@{}}Piensa de manera abstracta y formal en la resolu-\\ ción de problemas.\end{tabular} \\ \hline
SH06 & Saber Hacer & \begin{tabular}[c]{@{}l@{}}Sintetiza  distintos conceptos para abstraerse y\\ condensarse en un solo objeto matemático.\end{tabular} \\ \hline
SH07 & Saber Hacer & \begin{tabular}[c]{@{}l@{}}Visualiza y utiliza los conceptos geométricos-topo-\\ lógicos con el fin de encontrar nuevas aplicaciones\\ a la matemática y a la vida real.\end{tabular} \\ \hline
SH08 & Saber Hacer & \begin{tabular}[c]{@{}l@{}}Lee y entiende resultados escritos en lenguaje\\ matemático.\end{tabular} \\ \hline
SH09 & Saber Hacer & \begin{tabular}[c]{@{}l@{}}Plantea problemas reales en términos matemáticos,\\ mediante la modelación con miras a la aproxima-\\ ción, simulación y pronóstico.\end{tabular} \\ \hline
SH10 & Saber Hacer & \begin{tabular}[c]{@{}l@{}}Utiliza software científico para elaborar documentos,\\ crear algoritmos que le permitan plantear problemas,\\ obtener soluciones y divulgar los resultados.\end{tabular} \\ \hline
SH11 & Saber Hacer & \begin{tabular}[c]{@{}l@{}}Obtiene resultados mediante la ciencia de los datos\\ con el fin de encontrar relaciones e interpretaciones.\end{tabular} \\ \hline
SH12 & Saber Hacer & \begin{tabular}[c]{@{}l@{}}Expresa sus ideas en forma coherente mediante el\\ lenguaje matemático formal.\end{tabular} \\ \hline
SH13 & Saber Hacer & \begin{tabular}[c]{@{}l@{}}Usa la tecnología para explorar científicamente en \\ áreas del conocimiento matemático y afines.\end{tabular} \\ \hline
SH14 & Saber Hacer & \begin{tabular}[c]{@{}l@{}}Incorpora nuevos espacios de discusión científica y \\ matemática, o grupos de trabajo ya existentes.\end{tabular} \\ \hline
SH15 & Saber Hacer & \begin{tabular}[c]{@{}l@{}}Explora científicamente con el fin de entender temas\\ que no son de su especialidad.\end{tabular} \\ \hline
SH16 & Saber Hacer & \begin{tabular}[c]{@{}l@{}}Intercambia ideas con otros científicos y \\ matemáticos.\end{tabular} \\ \hline
SH17 & Saber Hacer & \begin{tabular}[c]{@{}l@{}}Actualiza sus conocimientos con el fin de mantener\\ vigencia en su área de especialidad.\end{tabular} \\ \hline
SH18 & Saber Hacer & \begin{tabular}[c]{@{}l@{}}Divulga y comunica el conocimiento matemático\\ que posee o las transformaciones que ha obtenido de\\ este a comunidades de especialistas o audiencias\\ más generales.\end{tabular} \\ \hline
SH19 & Saber Hacer & \begin{tabular}[c]{@{}l@{}}Relaciona la topología, el análisis y la geometría \\ con otras áreas de la Matemática y las ciencias en \\ general.\end{tabular} \\ \hline
SH20 & Saber Hacer & \begin{tabular}[c]{@{}l@{}}Relaciona el álgebra con otras áreas de la Matemá-\\ tica y las ciencias en general.\end{tabular} \\ \hline
SH21 & Saber Hacer & \begin{tabular}[c]{@{}l@{}}Visualiza los espacios euclídeos, hiperbólicos, \\ esféricos y protectivos, y las geometrías básicas en\\ dos, tres y más dimensiones.\end{tabular} \\ \hline
SH22 & Saber Hacer & \begin{tabular}[c]{@{}l@{}}Modela procesos reales (por ejemplo: fenómenos \\ aleatorios) mediante el uso de modelos matemáticos\\ tales como: procesos estocásticos, ecuaciones diferen-\\ ciales, métodos numéricos,etc.\end{tabular} \\ \hline
SH23 & Saber Hacer & \begin{tabular}[c]{@{}l@{}}Formula conjeturas con base en observaciones expe-\\ rimentales o teóricas sobre resultados plausibles.\end{tabular} \\ \hline
SS01 & Saber Ser & \begin{tabular}[c]{@{}l@{}}Es una persona que cuenta con la información y los\\ elementos necesarios para respetar la diversidad \\ humana independientemente de la etnia, género, \\ ideología, religión, estatus económico, orientación\\ sexual, nacionalidad o cualquier otra característica.\end{tabular} \\ \hline
SS02 & Saber Ser & \begin{tabular}[c]{@{}l@{}}Es una persona que escucha activamente para com-\\ prender la intención del mensaje que se le quiere\\ transmitir, y sabe comunicar de manera clara y \\ oportuna sus criterios y opiniones.\end{tabular} \\ \hline
SS03 & Saber Ser & \begin{tabular}[c]{@{}l@{}}Es una persona consciente de las diferencias que \\ existen entre los grupos sociales, así como de sus \\ necesidades y reconoce el valor de la organización y\\ de la participación ciudadana.\end{tabular} \\ \hline
SS04 & Saber Ser & \begin{tabular}[c]{@{}l@{}}Es una persona crítica y analítica que aplica los \\ principios éticos y de excelencia a las actividades\\ relacionadas con su disciplina.\end{tabular} \\ \hline
SS05 & Saber Ser & \begin{tabular}[c]{@{}l@{}}Es una persona que reconoce el valor del trabajo\\ interdisciplinario y que trabaja en equipo y de forma\\ colaborativa.\end{tabular} \\ \hline
SS06 & Saber Ser & \begin{tabular}[c]{@{}l@{}}Es una persona que concibe la realidad sociocultural\\ como una construcción fundada en la ciencia.\end{tabular} \\ \hline
SS07 & Saber Ser & \begin{tabular}[c]{@{}l@{}}Es una persona que concibe la ciencia y la matemática \\ como prácticas sociales interrelacionadas.\end{tabular} \\ \hline
SS08 & Saber Ser & \begin{tabular}[c]{@{}l@{}}Es una persona que respeta la diversidad matemática\\ y científica, y más generalmente, profesional y \\ cultural.\end{tabular} \\ \hline
SS09 & Saber Ser & \begin{tabular}[c]{@{}l@{}}Es una persona con disposición crítica hacia las diversas\\ prácticas científicas presentes en la sociedad.\end{tabular} \\ \hline
SS10 & Saber Ser & \begin{tabular}[c]{@{}l@{}}Es una persona que concibe la matemática como \\ fenómenos ligados a la cultura y la identidad de \\ una comunidad científica\end{tabular} \\ \hline
SS11 & Saber Ser & \begin{tabular}[c]{@{}l@{}}Es una persona que toma posiciones críticas ante las\\ instituciones académicas.\end{tabular} \\ \hline
SS12 & Saber Ser & \begin{tabular}[c]{@{}l@{}}Es una persona que tiene una visión integral de la \\ sociedad por medio del análisis de las producciones\\ científicas y matemáticas.\end{tabular} \\ \hline
SS13 & Saber Ser & \begin{tabular}[c]{@{}l@{}}Es una persona que cultiva y fomenta la creatividad \\ intelectual y científica a partir del conocimiento \\ matemático con el fin de ayudar a la comprensión de la \\ realidad circundante.\end{tabular} \\ \hline
SS14 & Saber Ser & \begin{tabular}[c]{@{}l@{}}Es una persona que conoce la importancia de la\\ matemática en la resolución de problemas propios y\\ prácticos del entorno.\end{tabular} \\ \hline
SS15 & Saber Ser & \begin{tabular}[c]{@{}l@{}}Es una persona que valora la importancia de conocer\\ la actualidad de la investigación científica y matemática.\end{tabular} \\ \hline
SS16 & Saber Ser & \begin{tabular}[c]{@{}l@{}}Es una persona que distingue las fuentes y personali-\\ dades académicas a las cuales acudir para solicitar \\ colaboración y ayuda en temas fuera de su área de \\ experticia.\end{tabular} \\ \hline
SS17 & Saber Ser & \begin{tabular}[c]{@{}l@{}}Es una persona que posee valor del autoaprendizaje y \\ que puede llegar a ser líder en la organización y dirección\\ de equipos de trabajo inter y multidisciplinario y dentro\\ de la matemática.\end{tabular} \\ \hline
\end{longtable}


\subsubsection{Presentacion de las “V” heurísticas de las identidades matemáticas presentes en la Escuela de Matematica de la Universidad de Costa Rica}

Guía de preguntas:

\begin{itemize}
    
    \color{blue} \item \color{black} ¿Cuáles son los conceptos fundamentales de esta área?

     \color{blue} \item \color{black}  ¿A partir de cuál acontecimiento histórico surgió la (se completa con la identidad)?
     
    \color{blue} \item \color{black} ¿Cómo se ubica usted a la (se completa con la identidad) dentro de la Matemática? 
    
    \color{blue} \item \color{black} ¿Cuáles áreas, ajenas a la (se completa con la identidad), han contribuido en el desarrollo de esta?
    
    \color{blue} \item \color{black} ¿A cuáles necesidades responde la investigación actual de la (se completa con la identidad)?
    
    \color{blue} \item \color{black} ¿Cuáles son las áreas más próximas en que la (se completa con la identidad) logra incidir?
   
    \end{itemize}

Según esta guía de preguntas, se extrajo la información y se presentó con el instrumento de una V heurística. La importancia de esta herramienta se observa en el transcurso del documento, donde se consideraron expresiones dadas por los docentes en las entrevistas o en los grupos focales respectivos.


\subsubsection{Tipos de unidades de aprendizaje}

\begin{longtable}{|
>{\columncolor[HTML]{8ED8F8}}l l|}
\hline
\multicolumn{2}{|c|}{\cellcolor[HTML]{00C0F3}\textbf{Tipos de Unidades de Aprendizaje}} \\ \hline
\multicolumn{1}{|c|}{\cellcolor[HTML]{8ED8F8}} & \begin{tabular}[c]{@{}l@{}}Actividad de aprendizaje en la cual los contenidos temáticos, \\ las actividades del estudiantado y el papel del profesorado se\\ orientan a la comprensión de cuerpos estructurados de \\ conocimientos, usualmente en forma de conceptos, teorías\\ y principios científicos.\end{tabular} \\ \cline{2-2} 
\multicolumn{1}{|c|}{\cellcolor[HTML]{8ED8F8}} & \begin{tabular}[c]{@{}l@{}}Los cursos teóricos se centran fundamentalmente en la \\ exposición, explicación, demostración y guía del docente\\ para contribuir a que sus estudiantes comprendan los \\ aspectos teóricos que se presentan.\end{tabular} \\ \cline{2-2} 
\multicolumn{1}{|c|}{\multirow{-3}{*}{\cellcolor[HTML]{8ED8F8}\textbf{Cursos teóricos}}} & \begin{tabular}[c]{@{}l@{}}Las actividades privilegian las horas de aula dirigidas por\\ un personal docente.\end{tabular} \\ \hline
\multicolumn{1}{|l|}{\cellcolor[HTML]{8ED8F8}} & \begin{tabular}[c]{@{}l@{}}Constituyen unidades de aprendizaje que se orientan a la \\ adquisición, por parte del estudiantado, de habilidades y \\ competencias propias de la aplicación del conocimiento, \\ así como a la solución de problemas, la creación de \\ artefactos, físicos o lógicos, la interpretación de situaciones\\ o la verificación de principios o leyes.\end{tabular} \\ \cline{2-2} 
\multicolumn{1}{|l|}{\cellcolor[HTML]{8ED8F8}} & \begin{tabular}[c]{@{}l@{}}Demandan del estudiantado participación activa, utilización\\ de equipos o ubicación de situaciones de aplicación práctica\\ de conocimientos y habilidades específicas.\end{tabular} \\ \cline{2-2} 
\multicolumn{1}{|l|}{\cellcolor[HTML]{8ED8F8}} & \begin{tabular}[c]{@{}l@{}}Demandan del profesorado la planificación de actividades, \\ la demostración y sobre todo la observación, correción y\\ guía de la actividad práctica que realiza el estudiantado.\end{tabular} \\ \cline{2-2} 
\multicolumn{1}{|l|}{\cellcolor[HTML]{8ED8F8}} & \begin{tabular}[c]{@{}l@{}}Demandan tiempo de ejercitación y aplicación de \\ procedimientos, en situaciones cercanas a la realidad\\ profesional.\end{tabular} \\ \cline{2-2} 
\multicolumn{1}{|l|}{\cellcolor[HTML]{8ED8F8}} & \begin{tabular}[c]{@{}l@{}}Demandan claridad de conceptos teóricos básicos y criterios\\ para discernir, en cada paso del proceso, que principios o \\ reglas se han de aplicar y por qué, para obtener determinados\\ resultados.\end{tabular} \\ \cline{2-2} 
\multicolumn{1}{|l|}{\cellcolor[HTML]{8ED8F8}} & \begin{tabular}[c]{@{}l@{}}Estos cursos favorecen la interacción entre estudiantes, la \\ flexibilidad en las horas de contacto entre docente y el \\ alumnado, según las necesidades durante el proceso.\end{tabular} \\ \cline{2-2} 
\multicolumn{1}{|l|}{\multirow{-7}{*}{\cellcolor[HTML]{8ED8F8}\textbf{Cursos prácticos}}} & \begin{tabular}[c]{@{}l@{}}A este grupo pertenecen los módulos de formación, los cursos\\ por proyectos, los laboratorios, los denominados talleres en\\ las áreas de artes o arquitectura, los proyectos colaborativos y,\\ tal vez, las denominadas prácticas profesionales, cuando \\ demandan una inserción del estudiantado en la comunidad\\ práctica real.\end{tabular} \\ \hline
\multicolumn{1}{|l|}{\cellcolor[HTML]{8ED8F8}} & \begin{tabular}[c]{@{}l@{}}Son actividades de aprendizaje de carácter práctico. En ellos, \\ el estudiante aplica conocimientos y practica habilidades y \\ destrezas mediante la realización de experimentos, análisis \\ de fenómenos, verificación de principios o modelos, cuya \\ realización requiere la utilización de instalaciones, \\ instrumentos y equipos especializados.\end{tabular} \\ \cline{2-2} 
\multicolumn{1}{|l|}{\cellcolor[HTML]{8ED8F8}} & \begin{tabular}[c]{@{}l@{}}En cuanto al desempeño docente, el profesorado que participa\\ en cursos laboratorio debe ser capaz de: planificar y disponer \\ previamente lo necesario para las actividades, demostrar un \\ uso adecuado de los equipos y materiales, contextualizar \\ la actividad en el ámbito de la profesión, guiar y facilitar\\ la ejecución de las actividades, analizar y evaluar el\\ desempeño del estudiantado.\end{tabular} \\ \cline{2-2} 
\multicolumn{1}{|l|}{\cellcolor[HTML]{8ED8F8}} & \begin{tabular}[c]{@{}l@{}}El profesorado utiliza como estrategias, formar, en el \\ estudiantado, el cumplimiento estricto de las normas \\ éticas en su propio comportamiento y la capacidad de\\ trabajar en equipo.\end{tabular} \\ \cline{2-2} 
\multicolumn{1}{|l|}{\multirow{-4}{*}{\cellcolor[HTML]{8ED8F8}\textbf{Cursos laboratorio}}} & \begin{tabular}[c]{@{}l@{}}También se comprende como clases de laboratorio a activi-\\ dades intensivas con medios de audio o video interactivos,\\ con la guía docente.\end{tabular} \\ \hline
\multicolumn{1}{|c|}{\cellcolor[HTML]{8ED8F8}} & \begin{tabular}[c]{@{}l@{}}El estudiantado realiza trabajos sobre una realidad concreta\\ y son sujetos u objetos propios de la acción profesional, en \\ lo cual es de particular importancia el proceso de solución o \\ ejecución, en el cual el estudiantado debe aplicar en forma\\ integrada conocimientos y destrezas, se sigue el principio \\ de aprender haciendo; y se privilegia el trabajo grupal, con \\ sistensis individuales.\end{tabular} \\ \cline{2-2} 
\multicolumn{1}{|c|}{\cellcolor[HTML]{8ED8F8}} & \begin{tabular}[c]{@{}l@{}}En este tipo de curso se requiere que el personal docente sea\\ capaz de: planificar e informar con claridad los objetivos de \\ aprendizaje, los procedimientos y los criterios para juzgar la\\ calidad del proceso y de los resultados; guiar el trabajo \\ individual y coordinar el trabajo grupal; dominar técnicas de\\ diálogo. preguntas y respuestas y de manejo de conflictos; \\ mantener apertura y flexibilidad para aceptar respuestas \\ divergentes y favorecer así la creatividad; ofrecer correcciones\\ oportunas e información de retorno permanente ajustada a los\\ requerimientos del estudiantado.\end{tabular} \\ \cline{2-2} 
\multicolumn{1}{|c|}{\multirow{-3}{*}{\cellcolor[HTML]{8ED8F8}\textbf{Curso taller}}} & \begin{tabular}[c]{@{}l@{}}El docente debe ser evaluado como planificador de las \\ actividades, incluyendo los aspectos logísticos que garanticen\\ una ejecución eficaz; su capacidad de interactuar con grupos \\ externos; su capacidad de solución de conflictos. Debe tener\\ habilidades para relacionarse con los centros de práctica.\end{tabular} \\ \hline
\multicolumn{1}{|l|}{\cellcolor[HTML]{8ED8F8}} & \begin{tabular}[c]{@{}l@{}}Es un curso de carácter práctico, de larga duración. Consiste\\ en la participación del estudiantado en una experiencia de \\ aplicación de su especialidad disciplinar. El estudiantado se\\ inserta en el ámbito profesional e interactúa con la realidad \\ social.\end{tabular} \\ \cline{2-2} 
\multicolumn{1}{|l|}{\cellcolor[HTML]{8ED8F8}} & \begin{tabular}[c]{@{}l@{}}Esta actividad le ofrece la oportunidad de investigar el contexto\\ y desarrollar competencias multidisciplinarias al integrar la\\ teoría y la práctica.\end{tabular} \\ \cline{2-2} 
\multicolumn{1}{|l|}{\cellcolor[HTML]{8ED8F8}} & \begin{tabular}[c]{@{}l@{}}Se espera que demuestre un alto nivel de compromiso con el \\ proceso, capacidad de priorizar y jerarquizar para la toma de \\ decisiones acertadas, y habilidad comunicativa para vincularse\\ con las instancias externas públicas y privadas.\end{tabular} \\ \cline{2-2} 
\multicolumn{1}{|l|}{\cellcolor[HTML]{8ED8F8}} & \begin{tabular}[c]{@{}l@{}}Para este tipo de curso se requiere que el docente y la docente\\ a cargo muestre estas competencias: ser mediador con el actor\\ social, establecer vínculos y facilitar las condiciones para que \\ se den las prácticas; tener habilidades para comunicarse \\ eficazmente, respetuoso con las instancias e instituciones que\\ prestan sus instalaciones; tener experiencia en el campo profe-\\ sional y disposición para realizar cambio e innovaciones; \\ asesorar y evaluar el desempeño del estudiantado; definir con\\ claridad cuáles son las intenciones de la experiencia, en qué\\ condiciones se realiza la práctica y qué responsabilidad tiene \\ cada uno.\end{tabular} \\ \cline{2-2} 
\multicolumn{1}{|l|}{\multirow{-5}{*}{\cellcolor[HTML]{8ED8F8}\textbf{Práctica profesional}}} & \begin{tabular}[c]{@{}l@{}}En la evaluación se hace necesario considerar también aspectos\\ relacionados con la cooperación entre el actor social y la propia \\ institución.\end{tabular} \\ \hline
\multicolumn{1}{|c|}{\cellcolor[HTML]{8ED8F8}} & \begin{tabular}[c]{@{}l@{}}Es una actividad de aprendizaje de carácter teórico. En ella el \\ estudiantado está orientado a fomentar el trabajo en equipo y \\ el aprendizaje autodirigido.\end{tabular} \\ \cline{2-2} 
\multicolumn{1}{|c|}{\cellcolor[HTML]{8ED8F8}} & \begin{tabular}[c]{@{}l@{}}Los participantes toman parte activa en el desarrollo del curso\\ con el profesor que dirige y coordina sobre la base de lecturas\\ críticas, investigaciones, bibliografía consultada y el conoci-\\ miento previo de cursos anteriores. Estos pueden ser temas \\ específicos de la disciplina o intereses particulares.\end{tabular} \\ \cline{2-2} 
\multicolumn{1}{|c|}{\cellcolor[HTML]{8ED8F8}} & \begin{tabular}[c]{@{}l@{}}Se espera desarrollar en el estudiantado el hábito de la lectura\\ con actitud crítica, fortalecer las habilidades de la escritura y \\ de la expresión oral, la discusión y el respeto en la opinión \\ fundamentada.\end{tabular} \\ \cline{2-2} 
\multicolumn{1}{|c|}{\multirow{-4}{*}{\cellcolor[HTML]{8ED8F8}\textbf{Seminario}}} & \begin{tabular}[c]{@{}l@{}}En este tipo de curso, se espera que cada docente sea capaz de:\\ corrdinar, orientar y evaluar la discusión y la exposición de los\\ alumnos; manejar técnicas de trabajo grupal para generar la \\ participación de los estudiantes; ser respetuoso con las diferen-\\ cias y la libertad ideológica; tener habilidades para investigar;\\ sintetizar los argumentos tratados en el aula, demostrando ser \\ un conocedor de la temática.\end{tabular} \\ \hline
\multicolumn{1}{|c|}{\cellcolor[HTML]{8ED8F8}} & \begin{tabular}[c]{@{}l@{}}Es un curso de tipo teórico-práctico. El curso tipo módulo\\ enfrenta al estudiantado a un conjunto de actividades de \\ aprendizaje alrededor de un eje temático, problema u objetivo,\\ que privilegia la práctica y la integración de conocimientos, lo\\ cual fortalece el carácter interdisciplinario y multidisciplinario\\ de la organización curricular.\end{tabular} \\ \cline{2-2} 
\multicolumn{1}{|c|}{\cellcolor[HTML]{8ED8F8}} & \begin{tabular}[c]{@{}l@{}}Permite al estudiantado fomentar el trabajo en equipo, el\\ espíritu de investigar e interpretar la realidad social, mejorar\\ la capacidad reflexiva, crítica y analítica.\end{tabular} \\ \cline{2-2} 
\multicolumn{1}{|c|}{\cellcolor[HTML]{8ED8F8}} & \begin{tabular}[c]{@{}l@{}}Desarrolla la capacidad de liderazgo, de autonomía y creatividad\\ bajo un comportamiento ético y responsable de sus actos.\end{tabular} \\ \cline{2-2} 
\multicolumn{1}{|c|}{\multirow{-4}{*}{\cellcolor[HTML]{8ED8F8}\textbf{Módulo}}} & \begin{tabular}[c]{@{}l@{}}Para este tipo de curso se espera que el profesorado sea capaz\\ de: comprometerse ética y socialmente con la tarea; manejar \\ metodologías participativas que orienten a los estudiantes al\\ abordaje temático; saber intervenir en un campo específico\\ ante problemas definidos; planificar e informar los objetivos\\ de aprendizaje, los procedimientos y los criterios para juzgar\\ la calidad del proceso y de los resultados; dar apoyo y \\ acompañamiento al trabajo individual y coordinar el manejo\\ de trabajo en equipo.\end{tabular} \\ \hline
\end{longtable}
